\subsection{无任意小有限子群的群的情形}

定理 2. --- 设 G 为满足下列条件的局部紧群：

\begin{enumerate}
\def\labelenumi{(\Alph{enumi})}
\setcounter{enumi}{11}
\tightlist
\item
  存在 \(e\) 在 \(G\) 中的一个邻域，它不含 \(G\) 的任何不约化为 \(e\) 的有限子群。
\end{enumerate}

于是下列性质成立：

\begin{enumerate}
\def\labelenumi{(\roman{enumi})}
\item
  G 的离散子群所成的集合 D 在 \(\Sigma\) 中是局部闭的（这等价于说它是局部紧的）。
\item
  为使 \(\mathbf{A}\) （\(\mathbf{D}\) 的闭子集）是紧的，必须且只需存在 \(\mathbf{U}\) （\(e\) 在 \(\mathbf{G}\) 中的一个邻域），使得 \(\mathbf{H} \cap \mathbf{U} = \{e\}\) 对每个子群 \(\mathbf{H} \in \mathbf{A}\) 都成立。
\item
  若此外 G 还由 e 的一个紧邻域生成，则由 G 中使 G/H 为紧的离散子群 H 所成的集合 \(D_{c}\) 在 \(\Sigma\) 中是局部闭的（从而是局部紧的）。
\end{enumerate}

我们有 \(\mathbf{D}_c = \mathbf{D} \cap \Sigma_c^0\) ，因此 (iii) 是 (i) 以及 No.~3 的 Th. 1 (ii) 的推论。

采用 No.~3 命题 7 的记号，为证 (i) 与 (ii)，只需证明：

命题 8. --- N 到 D 上的典范双射是一个同胚。

现在，若 \(\Gamma_d\) 是 G 的离散子群上的 Haar 测度所成的集合，则 D 典范地同胚于 \(\Gamma_d\) 中比 \(>0\) 的位似群的轨道空间 (GT, I, §5, No.~2, Prop. 4)。因此只需证明典范映射 \(\alpha \mapsto (\alpha(\{e\}), \alpha/\alpha(\{e\}))\) 是 \(\Gamma_d\) 到 \(\mathbf{R}_+^* \times \mathbf{N}\) 上的一个同胚，而这将由下面的引理得出：

引理 4. --- 若群 G 满足条件 (L)，则映射 \(\alpha \mapsto \alpha (\{e\})\) （\(\Gamma_d\) 到 \(\mathbf{R}_+^*\) 中的）是淡连续的。

考察一个测度 \(\alpha \in \Gamma_d\) ；设 \(V_0\) 是 \(e\) 在 \(G\) 中的一个相对紧开邻域，使得 \(H_\alpha \cap V_0 = \{e\}\) ，并且不含 \(G\) 的任何包含在 \(V_0\) 中的不约化为 \(e\) 的有限子群。设 \(V\) 是 \(e\) 的一个对称紧邻域，使得 \(V^3 \subset V_0\) ，又设 \(U\) 是 \(e\) 的一个对称邻域，使得 \(U^2 \subset V\) 。设 \(\varphi\) （相应地 \(\psi\) ）是 \(\mathcal{K}_+(G)\) 中的一个函数，取值于 \([0,1]\) ，在 \(V^3\) 上（相应地在点 \(e\) 处）等于 1，且支集包含在 \(V_0\) 中（相应地包含在 \(U\) 中）。\(\beta \in \Gamma_d\) 中使得 \(|\beta(\varphi) - \alpha(\varphi)| \leqslant \varepsilon\) 与 \(|\beta(\psi) - \alpha(\psi)| \leqslant \varepsilon\) 的测度所成的集合，是 \(W\) ，即 \(\alpha\) 的一个邻域。我们要证明：只要 \(\varepsilon\) 取得充分小，\(H_\beta \cap V = \{e\}\) 对每个 \(\beta \in W\) 都成立；由此将推出 \(\beta(\psi) = \beta(\{e\})\) ，从而 \(|\beta(\{e\}) - \alpha(\{e\})| \leqslant \varepsilon\) ，这就证明了引理。

只需证明：对 \(\beta \in \mathbf{W}\) ，

\[
(\mathrm{V} ^ {2} - \mathrm{V}) \cap \mathrm{H} _ {\beta} = \varnothing .\tag{8}
\]

因为，假设这一点已经确立：那么对 \(x\) 与 \(y\) （它们属于 \(\mathbf{V} \cap \mathbf{H}_{\beta}\) ），有 \(xy^{-1} \in \mathbf{V}^2 \cap \mathbf{H}_{\beta}\) ；但由 (8)，这蕴含 \(xy^{-1} \in \mathbf{V} \cap \mathbf{H}_{\beta}\) ；换言之，\(\mathbf{V} \cap \mathbf{H}_{\beta}\) 是 \(\mathbf{G}\) 的一个子群，它显然既离散又紧，因而是有限的；但这样一来，由 \(\mathbf{V}_0\) 的取法，这就蕴含确实有 \(\mathbf{V} \cap \mathbf{H}_{\beta} = \{e\}\) 。

我们用反证法，于是假设存在 \(z\) （\(\mathbf{V}^2 -\mathbf{V}\) 的一点）属于 \(\mathbf{H}_{\beta}\) ；由 U 与 V 的取法，在 G 中有 \(\psi (sz^{-1}) + \psi (s)\leqslant \varphi (s)\) ，因为关系式 \(z\notin \mathrm{U}^2\) 蕴含 \(\mathrm{U}z\cap \mathrm{U} = \emptyset\) 。由于

\[
\int \psi (s z ^ {- 1}) d \beta (s) = \int \psi (s) d \beta (s),
\]

由此推出 \(2\beta (\psi)\leqslant \beta (\varphi)\leqslant \alpha (\varphi) + \varepsilon\) ；但我们还有

\[
\beta (\psi) \geqslant \alpha (\psi) - \varepsilon ,
\]

而按构造有 \(\alpha(\varphi) = \alpha(\psi) = \alpha(\{e\})\) 。于是取 \(\varepsilon < \alpha(\{e\})/3\) 便得出矛盾。证毕。

形象地说，满足条件 (L) 的群 G 称为无任意小有限子群的群。*可以证明每个李群都满足条件 (L)；但这一条件并非李群的特征性质；例如，模 \(p\) 同余于 1 的 \(p\) 进整数所成的乘法群就满足 (L)。*

\subsection{交换群的情形}

设 G 为局部紧群，\(N \subset \Gamma^{0}\) 为 G 的离散子群上的规范化 Haar 测度所成的子空间，\(N_{c}\) 为 N 中对应于使 G/H 为紧的 G 的离散子群 H 的子集；于是在 No.~3 命题 6 的记号下有 \(N_{c} = N \cap \Gamma_{c}^{0}\) ；并且若群 G 由 e 的一个紧邻域生成，则由 No.~3 命题 6 推出，\(N_{c}\) 在 N 中是开的（从而由 No.~3 命题 7 的推论它是局部紧的），并且在 \(N_{c}\) 上映射 \(\alpha \mapsto \|\mu_{\alpha}\|\) 的限制是淡连续的。

命题 9. --- 设 G 为由 e 的一个紧邻域生成的局部紧交换群。为使 \(N_{c}\) 的子集 A 在 \(N_{c}\) 中相对紧，必须且只需它满足下列两个条件：

\begin{enumerate}
\def\labelenumi{(\roman{enumi})}
\item
  存在一个开邻域 \(\mathbf{U}\) （\(e\) 在 \(\mathbf{G}\) 中的），使得 \(\mathrm{H}_{\alpha} \cap \mathrm{U} = \{e\}\) 对所有 \(\alpha \in \mathbf{A}\) 都成立。
\item
  存在一个常数 \(k\) ，使得 \(\mu_{\alpha}(\mathrm{G} / \mathrm{H}_{\alpha}) \leqslant k\) 对所有 \(\alpha \in \mathbf{A}\) 都成立。若 \(\mathbf{A} \subset \mathbf{N}_c\) 在 \(\mathbf{N}_c\) 中相对紧，则它更在 \(\mathbf{N}\) 中相对紧，故条件 (i) 与 (ii) 的必要性由 No.~3 命题 6 与 7 得出（无需假设 \(\mathbf{G}\) 交换）。反之，设 \(\mathbf{A} \subset \mathbf{N}_c\) 满足这些条件；若 \(\overline{\mathbf{A}}\) 是 \(\mathbf{A}\) 在 \(\mathbf{N}\) 中的闭包，则由 No.~3 命题 7，\(\overline{\mathbf{A}}\) 是紧的；此外，由于 \(\alpha \mapsto \| \mu_{\alpha} \|\) 在 \(\Gamma^0\) 上关于淡拓扑是下半连续的 (No.~2, Prop. 4)，条件 (ii) 蕴含 \(\| \mu_{\alpha} \| \leqslant k\) 也对所有 \(\alpha \in \overline{\mathbf{A}}\) 成立。而由于 \(\mathbf{G}\) 交换，\(\mu_{\alpha} = \mu / \alpha\) 是群 \(G / H_{\alpha}\) 上的一个 Haar 测度，因此 \(G / H_{\alpha}\) 对每个 \(\alpha \in \overline{\mathbf{A}}\) 都是紧的 (Ch. VII, §1, No.~2, Prop. 2)。这就表明 \(\overline{\mathbf{A}} \subset N_c\) ，从而 \(\mathbf{A}\) 在 \(N_c\) 中相对紧。
\end{enumerate}

推论. --- 设 G 为由 e 的一个紧邻域生成且满足 No.~4 的条件 (L) 的局部紧交换群。

设 \(D_c\) 为使 G/H 为紧的 G 的离散子群 H 所成的集合，并且对每个 \(H \in D_c\) ，设 \(v(H)\) 为总质量 \(\mu_\alpha(G/H)\) ，其中 \(\mu_\alpha\) 是 \(\mu\) 关于 H 的规范化 Haar 测度 \(\alpha\) 的商测度。为使空间 \(D_c\) 的子集 A 在 \(D_c\) 中相对紧，必须且只需它满足下列两个条件：

\begin{enumerate}
\def\labelenumi{(\roman{enumi})}
\item
  存在 e 在 G 中的一个开邻域 U，使得 \(\mathbf{H} \cap \mathbf{U} = \{e\}\) 对所有 \(\mathbf{H} \in \mathbf{A}\) 都成立。
\item
  存在一个常数 \(k\) ，使得 \(v(\mathbf{H}) \leqslant k\) 对所有 \(\mathbf{H} \in \mathbf{A}\) 都成立。
\end{enumerate}

考虑到命题 9，这立即由下述事实得出：\(D_{c}\) 是 \(N_{c}\) 在 N 到 D 的典范双射下的像，并且在所作假设下该双射是一个同胚 (No.~4, Prop. 8)。

例. --- 取 \(G = R^{n}\) ，并取 \(\mu\) 为 Lebesgue 测度；命题 9 的推论的全部假设都得到满足。使 G/H 为紧的 G 的离散子群 H 正是秩为 n 的离散子群 (GT, VII, §1, No.~1, Th. 1)；这样的子群 H 由一个基 \((a_{i})_{1 \leqslant i \leqslant n}\) （\(R^{n}\) 的基）生成，并且

\[
v (\mathbf {H}) = | \det (a _ {1}, \dots , a _ {n}) |
\]

（行列式是关于 \(\mathbf{R}^n\) 的标准基取的）(Ch. VII, §2, No.~10, Th. 4)。空间 \(\mathbf{D}_c\) 在这里可按如下方式解释：每个子群 \(\mathbf{H} \in \mathbf{D}_c\) 都是 \(g \cdot \mathbf{Z}^n\) 型的子群（即子群 \(\mathbf{Z}^n\) 经某个元素 \(g \in \mathbf{GL}(n, \mathbf{R})\) 变换而得），而 \(\mathbf{GL}(n, \mathbf{R})\) 中使 \(\mathbf{Z}^n\) 保持稳定的子群可与 \(\mathbf{GL}(n, \mathbf{Z})\) 等同。因此，\(\mathbf{D}_c\) 作为（非拓扑的）齐性空间可以典范地与 \(\mathbf{GL}(n, \mathbf{R}) / \mathbf{GL}(n, \mathbf{Z})\) 等同。另一方面，\(\mathbf{GL}(n, \mathbf{R})\) 连续地作用在 \(\mathbf{R}^n\) 中，从而也作用于关于淡拓扑的 \(\mathcal{M}_{+}(\mathbf{R}^{n})\) 中 (§3, No.~3, Prop. 13)，因而作用于子空间 \(\mathrm{N}_c\) （\(\mathcal{M}_{+}(\mathbf{R}^{n})\) 的）中；此外，\(\mathrm{N}_c\) 到 \(\mathrm{D}_c\) 上的典范同胚 (No.~4, Prop. 8) 与 \(\mathbf{GL}(n, \mathbf{R})\) 的运算律相容。由于 \(\mathbf{GL}(n, \mathbf{R})\) 在无穷处可数且 \(\mathrm{D}_c\) 局部紧，上面定义的 \(\mathbf{GL}(n, \mathbf{R}) / \mathbf{GL}(n, \mathbf{Z})\) 到 \(\mathrm{D}_c\) 上的双射是一个同胚 (Ch. VII, App. I, Lemma 2)。于是命题 9 的推论给出了齐性空间 \(\mathbf{GL}(n, \mathbf{R}) / \mathbf{GL}(n, \mathbf{Z})\) 中紧性的一个判别法。

\subsection{闭子群空间拓扑的另一种解释}

设 F 为 G 的闭子集所成的集合；按如下方式在 F 上定义一个分离一致结构：对 G 的每个紧子集 K 以及 e 在 G 中的每个邻域 V，设 \(\mathrm{P}(\mathrm{K},\mathrm{V})\) 为 F 中同时满足下列两个条件的元素对 \((\mathrm{X},\mathrm{Y})\) 所成的集合：

\[
\mathrm{X} \cap \mathrm{K} \subset \mathrm{VY} \quad \text { and } \quad \mathrm{Y} \cap \mathrm{K} \subset \mathrm{VX}.\tag{9}
\]

我们来证明，诸 \(\mathrm{P}(\mathbf{K},\mathbf{V})\) 是一个分离一致结构 \(\mathcal{U}\) （在 \(\mathfrak{F}\) 上的）的邻偶基本系。GT, II, §1, No.~1 的公理 \((\mathrm{U}_{\mathrm{I}}^{\prime})\) 与 \((\mathrm{U}_{\mathrm{II}}^{\prime})\) 显然满足；此外，关系式 \(\mathbf{K} \subset \mathbf{K}'\) 与 \(\mathbf{V}' \subset \mathbf{V}\) 蕴含 \(\mathrm{P}(\mathbf{K}',\mathbf{V}') \subset \mathrm{P}(\mathbf{K},\mathbf{V})\) ；因此为验证 \((\mathrm{U}_{\mathrm{III}}^{\prime})\) ，可以只限于 \(\mathbf{V}\) 是 \(e\) 的对称紧邻域的情形，此时 VK 是紧的。设 \((\mathbf{X},\mathbf{Y}) \in \mathrm{P}(\mathrm{VK},\mathbf{V})\) 且 \((\mathbf{Y},\mathbf{Z}) \in \mathrm{P}(\mathrm{VK},\mathbf{V})\) ；那么 \(\mathbf{X} \cap \mathbf{K} \subset \mathbf{X} \cap \mathrm{VK} \subset \mathrm{VY}\) ，并且若 \(y \in \mathbf{Y}\) 使得 \(vy \in \mathbf{K}\) （对某个 \(v \in \mathbf{V}\) ），则必然 \(y \in \mathrm{VK}\) ，因此

\[
\mathrm{X} \cap \mathrm{K} \subset \mathrm{V} (\mathrm{Y} \cap \mathrm{VK});
\]

另一方面，\(Y \cap VK \subset VZ\) ，由此 \(X \cap K \subset V^{2}Z\) ，并且同理可证 \(Z \cap K \subset V^{2}X\) ，这就证明了 \((U_{III}^{\prime})\) 。最后，若 X、Y 是 F 的两个不同元素，则例如存在一点 \(a \in X\) 使得 \(a \notin Y\) ，从而存在 e 的一个对称紧邻域 V 使得 \(Va \cap Y = \varnothing\) ，即 \(a \notin VY\) ；更有 \((X,Y) \notin P(Va,V)\) ，这就完成了断言的证明。

这一点既已确立，我们在闭子群所成的集合 \(\Sigma\) （\(G\) 的）上考虑拓扑 \(\mathcal{T}\) ，它由刚刚定义的一致空间 \(\mathfrak{F}\) 的拓扑所诱导。我们将看到这一拓扑与 No.~3 中定义的拓扑相同。只需证明映射 \(\alpha \mapsto H_{\alpha}\) （\(\Gamma\) 到 \(\Sigma\) 中的映射）在 \(\Sigma\) 赋有拓扑 \(\mathcal{T}\) 时是连续的；因为这时该映射在 \(\Gamma_{\varphi}\) 上的限制（记号同 No.~1 命题 2）也是如此，而它是双射的；但由于 \(\Gamma_{\varphi}\) 紧且拓扑 \(\mathcal{T}\) 是分离的，映射 \(\alpha \mapsto H_{\alpha}\) （\(\Gamma_{\varphi}\) 到 \(\Sigma\) 中的）于是将是一个同胚。

于是设 \(\alpha_0\) 为 \(\Gamma\) 的一点，\(\Phi\) 为 \(\Gamma\) 上收敛于 \(\alpha_0\) 的一个滤子；我们要证明，关于 \(\Phi\) ，\(\mathrm{H}_{\alpha}\) 趋向于 \(\mathrm{H}_{\alpha_0}\) （对拓扑 \(\mathcal{T}\) 而言）。设 \(\mathbf{K}\) 为 \(\mathbf{G}\) 的紧子集，\(\mathbf{V}\) 为 \(e\) 在 \(\mathbf{G}\) 中的一个对称紧邻域；对每一点 \(x \in \mathrm{H}_{\alpha_0} \cap \mathrm{K}\) ，存在一个集合 \(\mathrm{M}(x) \in \Phi\) ，使得对每个 \(\alpha \in \mathrm{M}(x)\) 都有 \(\mathrm{V}x \cap \mathrm{H}_{\alpha} \neq \emptyset\) (No.~1, Lemma 2)，从而 \(\mathrm{V}x \subset \mathrm{V}^2\mathrm{H}_{\alpha}\) ；把 \(\mathrm{H}_{\alpha_0} \cap \mathrm{K}\) 用有限个集合 \(\mathrm{V}x_i\) 覆盖，就可以看出，若 \(\mathrm{M} = \bigcap \mathrm{M}(x_i)\) ，则 \(\mathrm{H}_{\alpha_0} \cap \mathrm{K} \subset \mathrm{V}^2\mathrm{H}_{\alpha}\) 对每个 \(\alpha \in \mathrm{M}\) 都成立。

反之，假设存在 e 在 G 中的一个开邻域 U，使得对每个集合 \(L \in \Phi\) ，至少有一个 \(\alpha \in L\) 满足 \(H_{\alpha} \cap K \not\subset UH_{\alpha_{0}}\) ；若 \(\omega(L)\) 是具有这一性质的 \(\alpha \in L\) 所成的集合，则诸 \(\omega(L)\) 构成滤子 \(\Phi'\) （\(\Gamma\) 上比 \(\Phi\) 更细的）的滤基，并且对每个 \(\alpha\) （它属于诸 \(\omega(L)\) （\(L \in \Phi\)）的并 E ），存在一个 \(t_{\alpha} \in H_{\alpha} \cap K\) ，它不属于 \(UH_{\alpha_0}\) ；对 \(\alpha \notin E\) ，取 \(t_{\alpha}\) 为 \(H_{\alpha}\) 的任一点。由于 \(K \cap \mathbb{C}(UH_{\alpha_0})\) 是紧的，将存在一个聚点 \(s\) （\(\alpha \mapsto t_{\alpha}\) 关于 \(\Phi'\) 的），它属于 \(K \cap \mathbb{C}(UH_{\alpha_0})\) ；但由于 \(\Phi'\) 收敛于 \(\alpha_0\) （在 \(\Gamma\) 中），这与 No.~1 的引理 3 矛盾。

\specialchapter{习题}

\exercisesection{}

\begin{enumerate}
\def\labelenumi{\arabic{enumi})}
\item
  设 \(\Gamma\) 是一个真 \(^{1}\) 闭凸锥（在 \(\mathbf{R}^{n}\) 中）。证明映射 \((x,y)\mapsto x+y\) （\(\Gamma\times\Gamma\) 到 \(\Gamma\) 中的）是逆紧的。由此推出，\(\Gamma\) 上任意两个测度关于映射 \((x,y)\mapsto x+y\) 都是可卷积的。
\item
  设 G 为局部紧群，\(\Gamma\) 为对 G 添加一个无穷远点 \(\omega\) 后得到的紧空间。把 G 的合成律扩充到 \(\Gamma\) ，办法是规定 \(x\omega = \omega x = \omega\) （对每个 \(x \in \Gamma\) ）。对每个测度 \(\mu\) （\(\Gamma\) 上的），一方面对应 G 上的一个有界测度 \(\mu_{1}\) ，另一方面对应一个复数 \(\mu(\{\omega\})\) 。证明：若用 \(*(\text{resp.}\widehat{*})\) 表示由 G（相应地 \(\Gamma\) ）中的乘法所定义的卷积，则 \((\mu\widehat{*}\nu)_{1} = \mu_{1}*\nu_{1}\) ，并且
\end{enumerate}

\[
(\mu \widehat {*} \nu) (\omega) = \mu (\omega) \nu_ {1} (G) + \nu (\omega) \mu_ {1} (G) + \mu (\omega) \nu (\omega)
\]

对任意两个测度 \(\mu\) 与 \(\nu\) （\(\Gamma\) 上的）成立。

\exercisesection{}

\begin{enumerate}
\def\labelenumi{\arabic{enumi})}
\tightlist
\item
  设 \((\mathbf{G}_{\iota})_{\iota \in \mathbf{I}}\) 是一族局部紧群，其中除有限个外都是紧的。设 \(U_{\iota}\) 是 \(\mathbf{G}_{\iota}\) 在一个局部凸空间 \(\mathbf{E}_{\iota}\) 中的连续线性表示。对每个 \(s = (s_{\iota}) \in \mathbf{G} = \prod_{\iota} \mathbf{G}_{\iota}\) ，设 \(U(s)\) 为自同态 \((x_{\iota}) \mapsto (U_{\iota}(s)x_{\iota})\) ，它定义在
\end{enumerate}

\(\mathbf{E} = \prod_{\iota}\mathbf{E}_{\iota}\) 上。证明 \(U\) 是 \(\mathbf{G}\) 在 \(\mathbf{E}\) 中的一个连续线性表示。设 \(\mathbf{E}'\) 为诸 \(E_{\iota}\) 的拓扑直和。设 \(V(s)\) 为 \(U(s)\) 在 \(E'\) 上的限制。证明 V 是 G 在 \(E'\) 中的一个连续线性表示。

\begin{enumerate}
\def\labelenumi{\arabic{enumi})}
\setcounter{enumi}{1}
\tightlist
\item
  设 \(U_{1}\) （相应地 \(U_{2}\) ）是局部紧群 \(G\) （相应地 \(H\) ）在一个局部凸空间 \(E\) （相应地 \(F\) ）中的连续线性表示。对 \(u \in \mathcal{L}(E; F)\) 、\(x \in G\) 、\(y \in H\) ，令
\end{enumerate}

\[
V (x, y) \cdot u = U _ {2} (y) \circ u \circ U _ {1} (x).
\]

证明映射 \((x,y)\mapsto V(x,y)\) 是群 \(G^{0}\times H\) 在赋有紧收敛拓扑的空间 \(\mathcal{L}(\mathrm{E};\mathrm{F})\) 中的一个连续线性表示。（利用 TVS, III, §4, No.~4 的 Prop. 9 以及如下事实：对 G 中紧的 K，\(U_{1}(K)\) 是等度连续的。）

¶ 3) 设 G 为局部紧群，U 是 G 在一个局部凸空间 E 中的连续线性表示，E' 是赋有强拓扑的 E 的对偶。

\begin{enumerate}
\def\labelenumi{\alph{enumi})}
\item
  证明对 G 的每个紧子集 K，\({}^{t}U(K)\) 是等度连续的。
\item
  设 \(\mathbf{F}\) 是由这样的 \(a' \in \mathbf{E}'\) 所成的集合：映射 \(s \mapsto {}^t U(s)a'\) （\(\mathbf{G}\) 到 \(\mathbf{E}'\) 中的）是连续的。证明 \(\mathbf{F}\) 是 \(\mathbf{E}'\) 的一个闭线性子空间，在 \({}^t U(\mathbf{G})\) 下稳定，并且限制到 \(\mathbf{F}\) 上、由 \(U\) 的逆步表示所得的表示是连续的。
\item
  假设 E 是拟完备的。设 \(\alpha\) 是 G 上的一个左 Haar 测度。证明 \(f \mapsto U(f \cdot \alpha)\) 是 \(\mathcal{K}(\mathrm{G})\) 到赋有有界收敛拓扑的 \(\mathcal{L}(\mathrm{E};\mathrm{E})\) 中的连续映射。（利用 Ch. VI, §1, No.~7 的 Prop. 17。）证明 F 在 \(E'\) 中弱稠密。（证明 \({}^{t}U(f)a' \in F\) 对所有 \(a' \in E'\) 与 \(f \in \mathcal{K}(\mathrm{G})\) 都成立，然后利用引理 4 的推论 3。）由此推出：若 E 是半自反的，则 U 在赋有强拓扑的 \(E'\) 中的逆步表示是连续的。
\item
  证明：若取 U 为 G 在 \(\mathrm{L}^{1}(\mathrm{G},\alpha)\) 中的左正则表示（\(\alpha\) 仍为 G 上的左 Haar 测度），则 F 是由一致连续函数所构成的 \(\mathrm{E}' = \mathrm{L}^{\infty}(\mathrm{G},\alpha)\) 的子空间。
\end{enumerate}

\begin{enumerate}
\def\labelenumi{\arabic{enumi})}
\setcounter{enumi}{3}
\item
  设 H 是一个 Hilbert 空间。称 G 在 H 中的连续表示 U 是酉的，如果自同态 \(U(s)\) 对所有 \(s \in G\) 都是酉的。对每个 \(\mu \in \mathcal{M}(G)\) ，设 \(\mu^{*}\) 为 \(\check{\mu}\) 的共轭测度。证明：若 \(\mu \in \mathcal{M}^{1}(G)\) ，则 \(U(\mu^{*}) = U(\mu)^{*}\) 。
\item
  设 \(\mathbf{G}\) 为局部紧群，\(\mathbf{H}\) 为 \(\mathbf{G}\) 的闭子群，\(U\) 为 \(\mathbf{H}\) 在一个局部凸空间 \(\mathbf{E}\) 中的连续线性表示。设 \(\mathbf{K}\) 为 \(\mathbf{G}\) 的一个紧子集。设 \(\mathcal{H}^U(\mathbf{K})\) 为 \(\mathbf{G}\) 上取值于 \(\mathbf{E}\) 、支集包含在 \(\mathrm{KH}\) 中且满足 \(f(xh) = U(h)^{-1}f(x)\) （ \(x \in \mathbf{G}, h \in \mathbf{H}\) ）的连续函数所成的空间。设 \(\mathcal{H}^U\) 为诸 \(\mathcal{H}^U(\mathbf{K})\) 的并，赋以在 \(\mathbf{K}\) 中一致收敛的诸拓扑（每个空间 \(\mathcal{H}^U(\mathbf{K})\) 上的）的归纳极限拓扑。对 \(f \in \mathcal{H}^U\) 与 \(s \in \mathbf{G}\) ，用下式定义 \(V(s)f \in \mathcal{H}^U\) ：
\end{enumerate}

\[
(V (s) f) (t) = f \left(s ^ {- 1} t\right).
\]

证明 \(V\) 是 \(\mathbf{G}\) 在 \(\mathcal{K}^U\) 中的一个连续线性表示。

\begin{enumerate}
\def\labelenumi{\arabic{enumi})}
\setcounter{enumi}{5}
\tightlist
\item
  设 \(\mathbf{G}\) 为局部紧群，\(\beta\) 为 \(\mathbf{G}\) 上一个相对不变的非零正测度，\(\chi\) 与 \(\chi'\) 是它的左、右乘子。对 \(f \in L_{\mathbf{C}}^{p}(\mathbf{G}, \beta)\) 与 \(s \in \mathbf{G}\) ，令
\end{enumerate}

\[
\begin{array}{c} (U (s) f) (x) = \chi (s) ^ {- 1 / p} f (s ^ {- 1} x) \\ (V (s) f) (x) = \chi^ {\prime} (s) ^ {- 1 / p} f (x s) \\ (S f) (x) = (\chi \chi^ {\prime}) (x) ^ {- 1 / p} \overline {{f (x ^ {- 1})}}. \end{array}
\]

则 \(U\) 与 \(V\) 是 \(\mathbf{G}\) 的线性表示，并且

\[
\begin{array}{c} S ^ {2} = 1, \| U (s) \| = \| V (s) \| = 1 \\ U (s) V (t) = V (t) U (s), S U (s) S = V (s) \end{array}
\]

对所有 \(s, t\) 属于 \(\mathbf{G}\) 成立。

\begin{enumerate}
\def\labelenumi{\arabic{enumi})}
\setcounter{enumi}{6}
\item
  设 E 是一个 Hilbert 空间，具有一个与 R 等势的标准正交基 \((e_{s})_{s\in\mathbf{R}}\) 。对每个 \(s\in R\) ，用 \(U(s)\) 表示 E 的等距变换，它满足 \(U(s)\cdot e_{t}=e_{s+t}\) （对所有 \(t\in R\) ）；R 在 E 中的线性表示 \(s\mapsto U(s)\) 不是连续的，尽管诸 \(U(s)\) 所成的集合是等度连续的。
\item
  设 \(G\) 为局部紧交换群，\(\mu\) 为 \(G\) 上的一个 Haar 测度，\(f\) 为一个 \(\mu\) -可测的有限数值函数（在 \(G\) 上）。假设对每个 \(s \in G\) ，数值函数
\end{enumerate}

\[
x \mapsto f (s x) - f (x)
\]

在 G 上连续。证明 f 这时是连续的。（用反证法；假设函数 f 在一点 \(x_{0} \in G\) 处不连续，先证明在 G 上存在一个以 e 为极限的滤子 F，使得

\[
\lim _ {\mathfrak {F}, s} | f (s x _ {0}) | = + \infty ,
\]

并由此推出 \(\lim_{\mathfrak{F},s}|f(sx)| = +\infty\) 对每个 \(x \in G\) 也成立。若 \(g = |f| / (1 + |f|)\) ，由最后这个结果推出它与下列事实矛盾：对每个紧子集 \(K \subset G\) ，

\[
\lim _ {\mathfrak {F}, s} \int_ {\mathrm{K}} | g (s x) - g (x) | d \mu (x) = 0.
\]

\begin{enumerate}
\def\labelenumi{\arabic{enumi})}
\setcounter{enumi}{8}
\item
  设 \(G\) 为局部紧群，\(\mu\) 为 \(G\) 上的一个左 Haar 测度，\(f\) 为一个 \(\mu\) -可积函数。设 \(\mathfrak{B}\) 为由 \(\mu\) -可积的、测度 \(> 0\) 的集组成、以 \(e\) 为极限的一个滤基。对每个 \(B \in \mathfrak{B}\) ，令 \(f_{\mathrm{B}}(t) = \frac{1}{\mu(\mathrm{B})} \int_{\mathrm{B}} f(st) d\mu(s)\) 。证明：对每个可积子集 \(A\) （\(G\) 的），有 \(\lim_{\mathfrak{B}} \int_{A} f_{\mathrm{B}}(t) d\mu(t) = \int_{A} f(t) d\mu(t)\) 。
\item
  设 \(\mathbf{G}\) 为局部紧群，\(\mathbf{E}\) 为分离的局部凸空间，\(\mathbf{E}'\) 为其对偶，\(U\) 为 \(\mathbf{G}\) 在 \(\mathbf{E}\) 中的一个线性表示，它对弱化拓扑 \(\sigma(\mathbf{E},\mathbf{E}')\) （\(\mathbf{E}\) 上的）是连续的。假设 \(\mathbf{E}\) 对 \(\sigma(\mathbf{E},\mathbf{E}')\) 是拟完备的，从而 \(U(\mu)\) 对每个测度 \(\mu \in \mathcal{C}'(\mathbf{G})\) 都有定义。证明双线性映射 \((\mu,x) \mapsto U(\mu) \cdot x\) 关于 \(\mathcal{C}'(\mathbf{G})\) 的等度连续子集是亚连续的。
\end{enumerate}

\exercisesection{}

\begin{enumerate}
\def\labelenumi{\arabic{enumi})}
\tightlist
\item
  在 \(\mathbf{R}^2\) 中，设 \(\lambda\) 与 \(\mu\) 为正测度
\end{enumerate}

\[
f \mapsto \int_ {0} ^ {+ \infty} f (x, 0) d x, \quad f \mapsto \int_ {0} ^ {+ \infty} f (- x, x) d x \quad \left(f \in \mathcal {K} (\mathbf {R} ^ {2})\right).
\]

证明 \(\lambda\) 与 \(\mu\) 是可卷积的。设 \(u\) 是同态 \((x,y) \mapsto x\) （\(\mathbf{R}^2\) 到 \(\mathbf{R}\) 上的）。证明 \(u\) 是 \(\lambda\) -逆紧的且是 \(\mu\) -逆紧的，但 \(u(\lambda)\) 与 \(u(\mu)\) 不是可卷积的。

\begin{enumerate}
\def\labelenumi{\arabic{enumi})}
\setcounter{enumi}{1}
\item
  设 X 是一个局部紧空间，一个局部紧群 G 在其中连续地从左边作用。设 E 是 \(\mathcal{M}(X)\) 的一个线性子空间，在诸 \(\boldsymbol{\gamma}(s)\) （ \(s \in G\) ）下稳定，赋有一个比 \(\mathcal{K}(X)\) 中紧收敛拓扑更细的拟完备局部凸拓扑。对每个 \(s \in G\) ，设 \(\boldsymbol{\gamma}_{\mathrm{E}}(s)\) 为 \(\boldsymbol{\gamma}(s)\) 在 E 上的限制。假设 \(\mu \in E\) 蕴含 \(|\mu| \in E\) ，且 G 在 E 中的表示 \(\gamma_{E}\) 是等度连续的。那么，若 \(\mu \in \mathbf{E}\) 且 \(\nu \in \mathcal{M}^1 (\mathbf{G})\) ，则 \(\nu\) 与 \(\mu\) 可卷积，并且 \(\nu * \mu = \gamma_{\mathbf{E}}(\nu)\mu \in \mathbf{E}\) 。（特别利用 Ch. VI, §1, No.~7 的 Prop. 17。）
\item
  对 \(x = (x_{1}, \ldots, x_{n}) \in \mathbf{R}^{n}\) ，置 \(|x|^{2} = x_{1}^{2} + \cdots + x_{n}^{2}\) 。设 \(M_{1}\) 为这样的测度 \(\mu\) （在 \(R^{n}\) 上）所成的集合：存在一个实数 k 使得函数 \((1 + |x|^{2})^{k}\) 是 \(\mu\) -可积的。设 \(M_{2}\) 为这样的测度 \(\nu\) （在 \(R^{n}\) 上）所成的集合：对每个 k，函数 \((1 + |x|^{2})^{k}\) 都是 \(\nu\) -可积的。证明：若 \(\mu \in M_{1}\) 且 \(\nu \in M_{2}\) ，则 \(\mu\) 与 \(\nu\) 可卷积，且 \(\mu * \nu \in M_{1}\) ；若 \(\mu \in M_{2}\) 且 \(\nu \in M_{2}\) ，则 \(\mu * \nu \in M_{2}\) 。（先证明：若 \(u \geqslant 0\) ，\(v \geqslant 0\) ，则
\end{enumerate}

\[
(1 + u ^ {2}) (1 + v ^ {2}) \geqslant \frac {1}{3} (1 + (u + v) ^ {2});
\]

由此推出，对任意的 \(x, y\) （在 \(\mathbf{R}^n\) 中），

\[
1 + | x | ^ {2} \leqslant 3 (1 + | y | ^ {2}) (1 + | x + y | ^ {2}).
\]

然后设 \(\mu \in \mathcal{M}_1\) ，\(\nu \in \mathcal{M}_2\) ，其中 \(\mu \geqslant 0\) ，\(\nu \geqslant 0\) 。设 \(f\) 是一个连续函数 \(\geqslant 0\) （在 \(\mathbf{R}^n\) 上）。存在一个 \(k\) 使得 \(\mu = (1 + |x|^2)^k \cdot \mu_1\) ，其中 \(\mu_1\) 有界；令 \(\nu_1 = (1 + |x|^2)^k \cdot \nu\) ，它是有界的。于是

\[
\int^ {*} f (x + y) d \mu (x) d \nu (y) \leqslant 3 ^ {k} \int^ {*} (1 + | x | ^ {2}) ^ {k} f (x) d (\mu_ {1} * \nu_ {1}) (x).
\]

由此得到 \(\mu\) 与 \(\nu\) 的可卷积性以及 \(\mu * \nu \in \mathcal{M}_1\) 这一事实。对 \(\mu, \nu\) （\(\mathcal{M}_2\) 中的）类似地论证。）

\begin{enumerate}
\def\labelenumi{\arabic{enumi})}
\setcounter{enumi}{3}
\item
  设 \(\mu\) 是 \(\mathbf{R}\) 上的 Lebesgue 测度，\(\nu\) 是 \([0, +\infty[\) 上的 Lebesgue 测度。设 \(x_1, x_2\) 属于 \(\mathbf{R}\) 。证明卷积 \(((\varepsilon_{x_1} - \varepsilon_{x_2}) * \nu) * \mu\) 有定义，但 \(\mu\) 与 \(\nu\) 不可卷积。证明卷积 \(\nu * ((\varepsilon_{x_1} - \varepsilon_{x_2}) * \mu)\) 与 \((\nu * (\varepsilon_{x_1} - \varepsilon_{x_2})) * \mu\) 都有定义，但当 \(x_1 \neq x_2\) 时它们不相等。
\item
  设 G 是紧群，\(\mu\) 是 G 上一个以 G 为支集的正测度，使得 \(\mu*\mu=\mu\) 。证明 \(\mu\) 是 G 的规范化 Haar 测度。（先证明 \(\|\mu\|=1\) 。然后，若 \(\mu\) 不是 G 的 Haar 测度，则存在一个函数 \(f\in\mathcal{K}_{+}(G)\) 使得
\end{enumerate}

\[
\int f (s) d \mu (s) \geqslant \int f (s t) d \mu (s)
\]

对所有 \(t \in \mathbf{G}\) 成立，且 \(\int f(s) d\mu(s) > \int f(st) d\mu(s)\) 对某个 \(t\) 成立。证明这时有 \((\mu * \mu)(f) > \mu(f)\) 。

\begin{enumerate}
\def\labelenumi{\arabic{enumi})}
\setcounter{enumi}{5}
\tightlist
\item
  \begin{enumerate}
  \def\labelenumii{\alph{enumii})}
  \tightlist
  \item
    设 \(I = [0,1]\) 。设 \(f\) 是一个连续函数 \(\geqslant 0\) （在 \(\mathbf{R}\) 上），其支集包含在 \([-1,0]\) 中。证明函数 \(\pmb{\gamma}(s)f|I\) （ \(s \in I\) ）所成的集合在 \(\mathcal{K}(I)\) 中有无穷秩。
  \end{enumerate}
\end{enumerate}

\begin{enumerate}
\def\labelenumi{\alph{enumi})}
\setcounter{enumi}{1}
\item
  设 \(f_1, \ldots, f_n\) 属于 \(\mathcal{K}(\mathbf{R})\) 。设 \(M\) 为这样的测度 \(\mu \in \mathcal{M}(I)\) 所成的集合：\(\mu(f_1) = \ldots = \mu(f_n) = 0\) 。证明函数 \(y \mapsto \int f(x + y) d\mu(x)\) （ \(y \in I\) ）（其中 \(\mu\) 取遍 \(M\) ）所成的集合在 \(\mathcal{K}(I)\) 中有无穷秩。（利用 \(a\) ）。
\item
  设 \(g_1, \ldots, g_p\) 属于 \(\mathcal{K}(\mathbf{R})\) 。由 \(b\) 推出，存在 \(\mu \in \mathbb{M}\) 与 \(\nu \in \mathcal{M}(\mathrm{I})\) 使得 \(\nu(g_1) = \ldots = \nu(g_p) = 0\) ，\((\nu * \mu)(f) \neq 0\) 。
\item
  由 c) 推出，映射 \((\mu,\nu)\mapsto\nu*\mu\) （\(\mathcal{M}(\mathbf{T})\times\mathcal{M}(\mathbf{T})\) 到 \(\mathcal{M}(\mathbf{T})\) 中的）不是淡连续的。
\end{enumerate}

\begin{enumerate}
\def\labelenumi{\arabic{enumi})}
\setcounter{enumi}{6}
\tightlist
\item
  设 \(G\) 为局部紧群。
\end{enumerate}

\begin{enumerate}
\def\labelenumi{\alph{enumi})}
\item
  证明：若正测度 \(\mu \in \mathcal{C}'(\mathrm{G})\) 在 \(\mathcal{C}'(\mathrm{G})\) 中有一个作为正测度的逆，则 \(\mu\) 必是点测度。
\item
  取 G 为有限群 \(\mathbf{Z}/2\mathbf{Z}\) ，给出一个非点测度 \(\mu \in \mathcal{M}(\mathbf{G})\) 的例子，它在 \(\mathcal{M}(\mathbf{G})\) 中可逆。
\end{enumerate}

\begin{enumerate}
\def\labelenumi{\arabic{enumi})}
\setcounter{enumi}{7}
\item
  设 G 为局部紧群；用 \(\mathcal{C}_{+}^{\prime}(G)\) 表示 G 上具有紧支集的正测度所成的集合。证明映射 \((\mu, \nu) \mapsto \mu * \nu\) （\(\mathcal{M}_{+}(G) \times \mathcal{C}_{+}^{\prime}(G)\) 到 \(\mathcal{M}_{+}(G)\) 中的）是连续的，其中 \(\mathcal{C}'(G)\) 赋有弱拓扑 \(\sigma(\mathcal{C}'(G), \mathcal{C}(G))\) 而 \(\mathcal{M}(G)\) 赋有淡拓扑 \(\sigma(\mathcal{M}(G), \mathcal{K}(G))\) 。（利用 Ch. III, §1 的 Exer. 10 以及 G 是仿紧的这一事实。）
\item
  \begin{enumerate}
  \def\labelenumii{\alph{enumii})}
  \tightlist
  \item
    设 G 为局部紧群，B 为 \(\mathcal{M}(\mathbf{G})\) 的一个有界子集，C 为 \(\mathcal{C}'(\mathbf{G})\) 的一个等度连续子集；证明：若 \(\mathcal{C}'(\mathbf{G})\) 赋有弱拓扑 \(\sigma(\mathcal{C}'(\mathbf{G}), \mathcal{C}(\mathbf{G}))\) ，且 \(\mathcal{M}(\mathbf{G})\) 赋有淡拓扑 \(\sigma(\mathcal{M}(\mathbf{G}), \mathcal{K}(\mathbf{G}))\) ，则映射 \((\mu, \nu) \mapsto \mu * \nu\) （B × C 到 \(\mathcal{M}(\mathbf{G})\) 中的）是连续的。（注意所有测度 \(\nu \in \mathbf{C}\) 的支集都在一个公共紧集内，并利用 Ch. III, §4, No.~3 中 Prop. 6 后的注。）
  \end{enumerate}
\end{enumerate}

\begin{enumerate}
\def\labelenumi{\alph{enumi})}
\setcounter{enumi}{1}
\tightlist
\item
  若在 \(\mathcal{M}^1 (\mathbf{R})\) 中令 \(\mu_n = \varepsilon_n\) ，\(\nu_{n} = \varepsilon_{-n}\) ，则序列 \((\mu_n)\) 与 \((\nu_{n})\) 对于弱拓扑 \(\sigma (\mathcal{M}^1 (\mathbf{R}),\overline{\mathcal{K}(\mathbf{R})})\) 淡地趋于 0，但序列 \((\mu_n*\nu_n)\) 对于淡拓扑不趋于 0。
\end{enumerate}

¶ 10) 对局部紧空间 T，用 \(\mathcal{C}^{\infty}(T)\) 表示 T 上有界连续数值函数所成的 Banach 空间。称 \(\mathcal{M}^{1}(T)\) 的子集 H 是挤压的，如果对每个 \(\varepsilon > 0\) ，存在 T 的一个紧子集 K，使得 \(|\mu|(T - K) \leqslant \varepsilon\) 对所有 \(\mu \in H\) 都成立。

\begin{enumerate}
\def\labelenumi{\alph{enumi})}
\item
  证明：若 \(\mathcal{M}^{1}(G)\) 的子集 H 是挤压的，且对 \(\mathcal{M}^{1}(T)\) 的范数所定义的拓扑有界，则它对拓扑 \(\sigma(\mathcal{M}^{1}(T),\mathcal{C}^{\infty}(T))\) 是相对紧的（注意 H 对 \(\sigma(\mathcal{M}^{1}(T),\mathcal{K}(T))\) 是相对紧的）。
\item
  再假设 T 是仿紧的，反之证明：若 H 是 \(\mathcal{M}^1 (\mathrm{T})\) 中对 \(\sigma (\mathcal{M}^1 (\mathrm{T}),\mathcal{C}^\infty (\mathrm{T}))\) 相对紧的子集，则 H 对 \(\mathcal{M}^1 (\mathrm{T})\) 的范数有界且是挤压的。（先考虑 \(\mathbf{T} = \mathbf{N}\) 的情形，应用 Ch. V, §5 的 Exer. 15。然后考察 T 在无穷处可数、是一列相对紧开集 \((\mathrm{U}_n)\) （满足 \(\overline{\mathrm{U}}_n\subset \mathrm{U}_{n + 1}\) ）之并的情形。用反证法证明，可以化归为这样的情形：对每个 \(n\) 存在一个定义在 T 上的连续数值函数 \(f_{n}\) ，其支集包含在 \(\mathrm{U}_{n + 1} - \overline{\mathrm{U}}_n\) 中，满足 \(\| f_n\| \leqslant 1\) ，并且存在一列属于 H 的测度 \((\mu_n)\) ，使得 \(\mu_{n}(f_{n})\geqslant a > 0\) 对所有 \(n\) 成立。于是考虑连续映射 \(u:\mathrm{L}^{1}(\mathbf{N})\to \mathcal{C}^{\infty}(\mathrm{T})\)
\end{enumerate}

使得 \(u((\xi_n)) = \sum_{n=0}^{\infty} \xi_n f_n\) ，并与对 \(T = N\) 已证明的结果相矛盾。

通过考虑 \(u\) 的转置。最后，当 \(T\) 是任意的仿紧局部紧空间时，\(T\) 是一族在无穷处可数的局部紧空间 \((T_{\alpha})\) 的拓扑和；对每个 \(\alpha\) ，令 \(m_{\alpha} = \sup_{\mu \in H} |\mu|(T_{\alpha})\) ；用反证法并利用前一种情形证明，必然有 \(m_{\alpha} = 0\) ，除对可数无穷多个指标 \(\alpha\) 外。）

\begin{enumerate}
\def\labelenumi{\alph{enumi})}
\setcounter{enumi}{2}
\tightlist
\item
  证明 \(b\) ) 的结论对 Ch. IV, §4 的 Exer. 16 h) 中定义的非仿紧局部紧空间不成立。
\end{enumerate}

¶ 11) 设 G 为局部紧群；在 \(\mathcal{M}^{1}(G)\) 上，用 \(\mathcal{T}_{\mathrm{I}}\) 表示拓扑 \(\sigma(\mathcal{M}^{1}(G), \overline{\mathcal{K}(G)})\) ，用 \(\mathcal{T}_{\mathrm{III}}\) 表示拓扑 \(\sigma(\mathcal{M}^{1}(G), \mathcal{C}^{\infty}(G))\) ，并用 \(\mathcal{M}_{\mathrm{I}}\) 、\(\mathcal{M}_{\mathrm{III}}\) 分别表示空间 \(\mathcal{M}^{1}(G)\) 赋有 \(\mathcal{T}_{\mathrm{I}}\) 与 \(\mathcal{T}_{\mathrm{III}}\) 后所得的空间。

\begin{enumerate}
\def\labelenumi{\alph{enumi})}
\item
  设 A 是 \(\mathcal{M}_{\mathrm{I}}\) 的一个有界子集，B 是 \(\mathcal{M}_{\mathrm{III}}\) 的一个相对紧子集。证明 \(\mathrm{A} \times \mathrm{B}\) 上映射 \((\mu, \nu) \mapsto \mu * \nu\) （\(\mathcal{M}_{\mathrm{I}} \times \mathcal{M}_{\mathrm{III}}\) 到 \(\mathcal{M}_{\mathrm{I}}\) 中的）的限制是连续的（利用 Exer. 10，以化归为估计积分 \(\iint f(st) d\mu(s) d\nu(t)\) ，当 \(f(st) = \sum_{i} u_{i}(s)v_{i}(t)\) 、其中 \(u_{i}\) 与 \(v_{i}\) 属于 \(\mathcal{K}(\mathrm{G})\) 时）。
\item
  给出一个 G 为紧（从而 \(T_{I} = T_{III}\) ）的例子，表明 \((\mu, \nu) \mapsto \mu * \nu\) 作为 \(M_{I} \times M_{III}\) 到 \(M_{I}\) 中的映射，既不关于 \(M_{I}\) 的相对紧子集亚连续，也不关于 \(M_{III}\) 的相对紧子集亚连续 (cf.~Exer. 6)。
\item
  设 \(A, B\) 为 \(\mathcal{M}_{\mathrm{III}}\) 的两个相对紧子集。证明 \(A \times B\) 上映射 \((\mu, \nu) \mapsto \mu * \nu\) （\(\mathcal{M}_{\mathrm{III}} \times \mathcal{M}_{\mathrm{III}}\) 到 \(\mathcal{M}_{\mathrm{III}}\) 中的）的限制是连续的（方法同 a)）。
\item
  用 E 表示由 G 的开子集的特征函数的线性组合构成的 \(\mathbf{R}^{\mathrm{G}}\) 的子空间，用 \(\mathcal{T}_{\mathrm{IV}}\) 表示拓扑 \(\sigma(\mathcal{M}^{1}(\mathrm{G}),\mathrm{E})\) ，并用 \(\mathcal{M}_{\mathrm{IV}}\) 表示空间 \(\mathcal{M}^{1}(\mathrm{G})\) 赋有 \(\mathcal{T}_{\mathrm{IV}}\) 后所得的空间。回忆一下，\(\mathcal{T}_{\mathrm{III}}\) 与 \(\mathcal{T}_{\mathrm{IV}}\) 在由正测度组成的 \(\mathcal{M}^{1}(\mathrm{G})\) 的有界子集上诱导的拓扑一般是不相同的 (Ch. V, §5, Exer. 16 c))。再回忆一下，\(\mathcal{M}_{\mathrm{IV}}\) 的紧子集与 Banach 空间 \(\mathcal{M}^{1}(\mathrm{G})\) 赋弱化拓扑 \(\sigma(\mathcal{M}^{1}(\mathrm{G}),(\mathcal{M}^{1}(\mathrm{G}))'\) 时 \(\mathcal{M}^{1}(\mathrm{G})\) 的紧子集相同 (Ch. VI, §2, Exer. 12)。证明：若 A、B 是 \(\mathcal{M}_{\mathrm{IV}}\) 的两个相对紧子集，则映射 \((\mu,\nu)\mapsto\mu*\nu\) （\(\mathcal{M}_{\mathrm{IV}}\times\mathcal{M}_{\mathrm{IV}}\) 到 \(\mathcal{M}_{\mathrm{IV}}\) 中的）在 A × B 上的限制是连续的。（只限于 A 与 B 是紧的情形，并且先证明这时 A × B 在上述映射下的像对 \(\sigma(\mathcal{M}^{1}(\mathrm{G}),(\mathcal{M}^{1}(\mathrm{G}))'\) 是紧的；为此，应用 Eberlein 定理与 Šmulian 定理 (TVS, IV, §5, No.~3)，从而化归为证明：若 \((\mu_{n})\) 、\((\nu_{n})\) 是在 \(\mathcal{M}_{\mathrm{IV}}\) 中收敛于 0 的两个序列，则序列 \((\mu_{n}*\nu_{n})\) 亦然；利用 Ch. V, §5 的 Prop. 12 与 Exer. 15。最后，为证 \((\mu,\nu)\mapsto\mu*\nu\) 对 \(\mathcal{T}_{\mathrm{IV}}\) 连续，利用 c) 以及 \(\mathcal{T}_{\mathrm{III}}\) 比 \(\mathcal{T}_{\mathrm{IV}}\) 粗这一事实。）
\item
  取 \(\mathbf{G} = \mathbf{R}^2\) ；设 \(a, b\) 是 \(\mathbf{G}\) 在 \(\mathbf{R}\) 上的标准基的向量，\(\rho_n\) 是区间 \(\mathbf{I} = [0, \pi]\) （\(\mathbf{R}\) 中的）上以函数 \(\sin(2^n x)\) 为密度（关于 Lebesgue 测度）的测度，\(\mu_n\) 是测度 \(\rho_n \otimes \varepsilon_0\) （\(\mathbf{R} \times \mathbf{R}\) 上的）；另一方面，设 \(\nu_n = \varepsilon_{b/2^n} - \varepsilon_0\) （\(\mathbf{G}\) 上的）。证明序列 \((\mu_n)\) 在 \(\mathcal{M}_{\mathrm{IV}}\) 中趋于 0，序列 \((\nu_n)\) 在 \(\mathcal{M}_{\mathrm{III}}\) 中趋于 0，但序列 \((\mu_n * \nu_n)\) 在 \(\mathcal{M}_{\mathrm{IV}}\) 中不趋于 0 (cf.~Ch. V, §5, Exer. 16)。
\item
  给出一个 G 为紧且 \((\mu,\nu)\mapsto\mu*\nu\) 作为 \(M_{IV}\times M_{IV}\) 到 \(M_{IV}\) 中的映射不关于 \(M_{IV}\) 的紧子集亚连续的例子（方法同 Exer. 6，注意对给定的 \(f\in E\) ，存在紧子集 \(H\subset M_{IV}\) 使得函数 \(t\mapsto\int f(st)d\mu(s)\) （其中 \(\mu\) 取遍 H）所成的集合在 R 上有有限秩）。
\end{enumerate}

\begin{enumerate}
\def\labelenumi{\arabic{enumi})}
\setcounter{enumi}{11}
\tightlist
\item
  设 \(G\) 为一个非幺模的局部紧群。
\end{enumerate}

\begin{enumerate}
\def\labelenumi{\alph{enumi})}
\item
  证明存在一个有界正测度 \(\mu\) （\(\mathbf{G}\) 上的），使得 \(\Delta_{\mathbf{G}} \cdot \mu\) 不是有界的（取 \(\mu\) 为离散的）。
\item
  设 \(\mu'\) 是 G 上的一个左 Haar 测度。那么 \(\mu\) 与 \(\mu'\) 可卷积 (Prop. 5)。证明 \(\mu'\) 与 \(\mu\) 不可卷积。
\end{enumerate}

\begin{enumerate}
\def\labelenumi{\arabic{enumi})}
\setcounter{enumi}{12}
\tightlist
\item
  设 \(r\) 是一个满足 \(0 < r < 1\) 的数；对每个整数 \(n \geqslant 1\) ，用 \(\lambda_{n,r}\) 表示测度 \((\varepsilon_{r^n} + \varepsilon_{-r^n}) / 2\) （\(\mathbf{R}\) 上的），并令 \(\mu_{n,r} = \lambda_{1,r} * \lambda_{2,r} * \dots * \lambda_{n,r}\) 。
\end{enumerate}

\begin{enumerate}
\def\labelenumi{\alph{enumi})}
\item
  证明序列 \((\mu_{n,r})\) 淡收敛到 R 上的一个测度 \(\mu_{r}\) ，其支集包含在 I = {[}-1,+1{]} 中（证明对 R 的每个区间 U，序列 \((\mu_{n,r}(U))\) 收敛）。
\item
  证明当 \(r < 1/2\) 时，测度 \(\mu_r\) 与 \(\mathbf{R}\) 上的 Lebesgue 测度互奇异，但 \(\mu_{1/2}\) 是 Lebesgue 测度在 I 上诱导的测度。
\item
  设 \(\nu_{1/4}\) 是 \(\mu_{1/4}\) 在位似 \(t \mapsto 2t\) （\(\mathbf{R}\) 中的）下的像。证明 \(\mu_{1/4} * \nu_{1/4} = \mu_{1/2}\) ，尽管 \(\mu_{1/4}\) 与 \(\nu_{1/4}\) 都与 Lebesgue 测度互奇异（利用 Exer. 11 a)）。
\end{enumerate}

\exercisesection{}

¶ 1) 设 G 为局部紧群，β 为 G 上的一个左 Haar 测度，A 为 G 的一个 β-可测子集，ν 为 G 上一个非零正测度。假设对每个 \(s \in G\) ，sA 都是局部 ν-可忽略的。证明 A 是局部 β-可忽略的。（化归为 A 相对紧且 ν 有界的情形。证明 ν 与 \(\varphi_{A} \cdot \beta\) 可卷积，且 \(\nu *^{\beta} \varphi_{A} = 0\) ，从而 \(0 = \|\nu * \varphi_{A} \cdot \beta\| = \|\nu\| \cdot \|\varphi_{A} \cdot \beta\|\) 。）在承认连续统假设的前提下证明：若不假设 A 是 β-可测的，这一结果可能不成立。（取 \(G = R^{2}\) ，对 ν 取子群 \(R \times \{0\}\) 上的 Haar 测度，并应用 Ch. V, §8 的 Exer. 7 c)。）

\begin{enumerate}
\def\labelenumi{\arabic{enumi})}
\setcounter{enumi}{1}
\tightlist
\item
  设 H 是赋有离散拓扑的加法群 R。设 G 为局部紧群 \(R \times R \times H\) 。设 \(\alpha\) 是 G 上的一个 Haar 测度，\(\beta\) 是 \(R \times H\) 上的一个 Haar 测度，而 \(\mu = \varepsilon_{0} \otimes \beta \in \mathcal{M}(G)\) 。在 G 上构造一个函数 \(f \geqslant 0\) ，它局部 \(\alpha\) -可忽略（从而 \(\mu\) 与 f 可卷积），但 f 的任何平移都不是 \(\mu\) -可测的。（模仿 Ch. V, §3, Exer. 4 的构造。）
\end{enumerate}

¶ 3) 设 G 为局部紧群。

\begin{enumerate}
\def\labelenumi{\alph{enumi})}
\tightlist
\item
  设 \(\mu\) 是 \(\mathbf{G}\) 上一个非零的有界正测度，满足 \(\mu * \mu = \mu\) 。证明支集 \(\mathbf{S}\) （\(\mu\) 的）是紧的。（取 \(f \in \mathcal{K}_{+}(\mathbf{G})\) 且 \(f \neq 0\) ；在 \(\mathbf{G}\) 上选定一个左 Haar 测度，卷积都将关于它取；我们有 \(\mu * f \in \overline{\mathcal{K}(\mathbf{G})}\) ；设 \(x \in \mathbf{S}\) 使得 \((\mu * f)(x) = \sup_{y \in \mathbf{S}} (\mu * f)(y)\) ；证明
\end{enumerate}

\((\mu * f)(y) = (\mu * f)(x)\) 对每个 \(y \in S\) 成立。

\begin{enumerate}
\def\labelenumi{\alph{enumi})}
\setcounter{enumi}{1}
\tightlist
\item
  证明 S 是 G 的紧子群，且 \(\mu\) 是 S 的规范化 Haar 测度。（利用 a)、GT, III, §4 的 Exer. 21 以及 §3 的 Exer. 5。）
\end{enumerate}

\begin{enumerate}
\def\labelenumi{\arabic{enumi})}
\setcounter{enumi}{3}
\tightlist
\item
  设 \(G\) 为局部紧群。对每个 \(t \in \mathbf{R}_+^*\) ，设 \(\mu_t\) 是 \(G\) 上一个非零的有界正测度。假设映射 \(t \mapsto \mu_t\) 对拓扑 \(\sigma(\mathcal{M}^1(G), \overline{\mathcal{K}(G)})\) 连续，且 \(\mu_{s+t} = \mu_s * \mu_t\) （ \(s, t\) 属于 \(\mathbf{R}_+^*\) ）。
\end{enumerate}

\begin{enumerate}
\def\labelenumi{\alph{enumi})}
\item
  证明存在一个数 \(c \in \mathbf{R}\) 使得 \(\| \mu_t \| = \exp(ct)\) 。（注意 \(t \mapsto \| \mu_t \|\) 是下半连续的，且 \(\| \mu_{s+t} \| = \| \mu_s \| \cdot \| \mu_t \|\) 。利用 Prop. 18。）
\item
  假设 \(c = 0\) 。证明 \(\mu_t\) 当 \(t \to 0\) 时弱收敛到 \(G\) 的某个紧子群的规范化 Haar 测度。（利用 \(\mathcal{M}^1(G)\) 的单位球的弱紧性，并对每个弱聚点 \(\mu\) （\(t \mapsto \mu_t\) 关于 \(R_+^*\) 中 0 的邻域滤子的），证明 \(\mu_t * \mu = \mu_t\) 对每个 \(t\) 成立，然后 \(\mu * \mu = \mu\) ；接着应用 Exer. 3。）
\end{enumerate}

¶ 5) 设 G 为在无穷处可数的局部紧群，β 为 G 上的一个左 Haar 测度，\(a \in R_{+}^{*}\) ，且 \((\nu_{u})_{0 \leqslant u \leqslant a}\) 为 G 上满足下列条件的一族正测度：

\begin{enumerate}
\def\labelenumi{(\roman{enumi})}
\item
  \(\nu_0 = \varepsilon_e\) ；\(\nu_{u + v} = \nu_u * \nu_v\) （若 \(u + v \leqslant a\) ）；\(\nu_u = \stackrel{\vee}{\nu}_u\) ；
\item
  对 \(0 < u \leqslant a\) ，\(\nu_{u} = f_{u} \cdot \beta\) ，其中 \(f_{u} \geqslant 0\) 是下半连续的；
\item
  \(\nu_{u}\) 是 \(u\) 的一个淡连续函数。
\end{enumerate}

\begin{enumerate}
\def\labelenumi{\alph{enumi})}
\tightlist
\item
  设 \(f \in \mathcal{K}_{+}(\mathrm{G})\) 。证明对 \(0 < u \leqslant a\) ，
\end{enumerate}

\[
\int \left(f _ {u / 2} * f\right) (z) ^ {2} d \beta (z) = \int f (x) \left(f _ {u} * f\right) (x) d \beta (x).\tag{1}
\]

\begin{enumerate}
\def\labelenumi{\alph{enumi})}
\setcounter{enumi}{1}
\item
  设 \(f \in \mathcal{K}(\mathbf{G})\) 。证明 \(f_{u/2} * f\) 有 \(\beta\) -可积的平方（对 \(0 < u \leqslant a\) ），且 (1) 仍成立。
\item
  设 \(f \in \mathcal{H}(\mathrm{G})\) 满足 \(\nu_{a} * f = 0\) 。证明 \(f = 0\) 。（有 \(\nu_{a/2^n} * f = 0\) 对每个整数 \(n > 0\) （由 \(b\) )），故由 §3, No.~3 的注 1 得 \(f = 0\) 。）
\item
  设 \(\nu \in \mathcal{C}'(\mathrm{G})\) 满足 \(\nu_{a} * \nu = 0\) 。证明 \(\nu = 0\) 。（正则化 \(\nu\) （用 \(\mathcal{H}(\mathrm{G})\) 中的函数）并应用 \(c\) ）。）
\end{enumerate}

\begin{enumerate}
\def\labelenumi{\arabic{enumi})}
\setcounter{enumi}{5}
\item
  设 G 为局部紧群。证明：若 \(\mathcal{K}(\mathrm{G})\) 对卷积是交换的，则 G 是交换群。（用正则化证明 \(\mathcal{C}'(\mathrm{G})\) 是交换的，并把这一点应用到测度 \(\varepsilon_{s}\) 上，其中 \(s \in G\) 。）
\item
  设 \(G\) 为局部紧群，\(\beta\) 为 \(G\) 上的一个左 Haar 测度。证明：代数 \(L^1(G, \beta)\) 有单位元当且仅当 \(G\) 是离散的。（设 \(G\) 非离散，并取 \(f_0 \in L^1(G, \beta)\) 。存在一个紧邻域 \(V\) （\(e\) 的）使得
\end{enumerate}

\[
\int_ {V} | f _ {0} (x) | d \beta (x) <   1.
\]

设 \(\mathbf{U}\) 是 \(e\) 的一个对称紧邻域，使得 \(\mathbf{U}^2 \subset \mathbf{V}\) 。那么对几乎每个 \(x \in \mathbf{U}\) ，

\[
\left| \left(\varphi_ {\mathrm{U}} * f _ {0}\right) (x) \right| = \int_ {\mathrm{U}} \left| f _ {0} \left(y ^ {- 1} x\right) \right| d \beta (y) \leqslant \int_ {\mathrm{V}} \left| f _ {0} (x) \right| d \beta (x) <   1,
\]

因而 \(f_0\) 不是 \(L^1(G, \beta)\) 的单位元。

\begin{enumerate}
\def\labelenumi{\arabic{enumi})}
\setcounter{enumi}{7}
\tightlist
\item
  设 G 为在无穷处可数的局部紧群，连续地从左边作用在一个波兰局部紧空间 T 中。设 \(\nu\) 是 T 上一个在 G 下拟不变的正测度。设 R 是 T 上一个 \(\nu\) -可测的等价关系，与 G 相容。于是存在 (Ch. VI, §3, No.~4, Prop. 2) 一个波兰局部紧空间 B 和一个 T 到 B 中的 \(\nu\) -可测映射 p，使得 \(R\{x,y\}\) 等价于 \(p(x)=p(y)\) 。设 \(\nu'\) 是 \(\nu\) 在 p 下的伪象测度，并设 \(b\mapsto\lambda_{b}\) （ \(b\in B\) ）是 \(\nu\) 关于 R 的一个分解。证明 \(\lambda_{b}\) 对几乎每个 \(b\in B\) 在 G 下拟不变。
\end{enumerate}

（设 \(\chi\) 是 \(\mathbf{G} \times \mathbf{T}\) 上满足 Ch. VII, §1 的 Exer. 13 的条件的函数。证明对每个 \(s \in \mathbf{G}\) ，存在 B 的一个 \(\nu'\) -可忽略子集 \(\mathbf{N}(s)\) ，使得 \(\chi(s^{-1}, \cdot)\) 是局部 \(\lambda_b\) -可积的（对 \(b \notin \mathbf{N}(s)\) ）；为此，注意对每个 \(\psi \in \mathcal{K}(\mathbf{T})\) ，函数 \(x \mapsto \psi(x)\chi(s^{-1}, x)\) 是 \(\lambda_b\) -可积的，除了 \(b\) 属于某个 \(\nu'\) -可忽略集 \(\mathbf{N}(s, \psi)\) 以外，并利用 Ch. VI, §3, No.~1 的引理 1。令 \(\lambda_{b,s}' = 0\) （若 \(b \in \mathbf{N}(s)\) ），并令 \(\lambda_{b,s}' = \chi(s^{-1}, \cdot) \cdot \lambda_b\) （若 \(b \notin \mathbf{N}(s)\) ）。证明映射 \(b \mapsto \lambda_{b,s}'\) 是 \(\nu'\) -适配的（利用 Ch. VI, §3, No.~1 的引理 3），且 \(\gamma(s)\beta = \int \lambda_{b,s}' d\nu'(b)\) 。另一方面证明 \(\gamma(s)\beta = \int \gamma(s)\lambda_b d\nu'(b)\) ，并由此推出：对每个 \(s \in \mathbf{G}\) ，有 \(\gamma(s)\lambda_b = \chi(s^{-1}, \cdot) \cdot \lambda_b\) 对几乎每个 \(b\) ，从而对几乎每个 \(b\) 有 \(\gamma(s)\lambda_b = \chi(s^{-1}, \cdot) \cdot \lambda_b\) 对几乎每个 \(s\) 。然后利用 §4 的 Prop. 17 的推论 2 推断：对几乎每个 \(b\) ，\(\gamma(s)\lambda_b\) 都与 \(\lambda_b\) 等价（对所有 \(s \in \mathbf{G}\) ）。）

证明：若 \(\nu\) 在 G 下以乘子 \(\chi\) 相对不变，则对几乎每个 b，\(\lambda_{b}\) 以乘子 \(\chi\) 相对不变。

\begin{enumerate}
\def\labelenumi{\arabic{enumi})}
\setcounter{enumi}{8}
\item
  设 \(G\) 为局部紧群，\(\beta\) 为 \(G\) 上的一个左 Haar 测度，\(A\) 与 \(B\) 为两个 \(\beta\) -可积集，满足 \(\beta(A) \leqslant \beta(B)\) 且 \(\beta^{*}(A^{-1}) < +\infty\) 。证明存在
\item
  不相交集合 \(N, K_{1}, K_{2}, \ldots\) 覆盖 \(A\) ，其中 \(N \beta\) -可忽略而诸 \(K_{n}\) 是紧的；
\item
  覆盖 B 的不相交集合 \(\mathbf{N}^{\prime}\) 、\(\mathbf{K}_1^{\prime}\) 、\(\mathbf{K}_2^{\prime}\) 、\ldots{} ，其中诸 \(\mathbf{K}_n^{\prime}\) 是紧的；\\
\item
  元素 \(s_n \in G\) 使得 \(\mathbf{K}_n^{\prime} = s_n\mathbf{K}_n\) 。
\end{enumerate}

（利用 \(\beta (x\mathbf{A}\cap \mathbf{B})\) 连续地依赖于 \(x\) 这一事实以及

\[
\int \beta (x \mathrm{A} \cap \mathrm{B}) d \beta (x) = \beta (\mathrm{A} ^ {- 1}) \beta (\mathrm{B}),
\]

证明：若 \(\beta (\mathbf{A})\neq 0\) ，则存在 \(x\in \mathbf{G}\) 使得 \(\beta (x\mathbf{A}\cap \mathbf{B})\neq 0\) 。）

\begin{enumerate}
\def\labelenumi{\arabic{enumi})}
\setcounter{enumi}{9}
\tightlist
\item
  设 \(\mathbf{G}\) 为局部紧群，\(\beta\) 为 \(\mathbf{G}\) 上的一个左 Haar 测度。对任意两个 \(\beta\) -可积集 \(\mathbf{A}, \mathbf{B}\) ，令 \(\rho(\mathbf{A}, \mathbf{B}) = \beta(\mathbf{A} \cup \mathbf{B}) - \beta(\mathbf{A} \cap \mathbf{B})\) 。
\end{enumerate}

\begin{enumerate}
\def\labelenumi{\alph{enumi})}
\item
  设 A 是一个 \(\beta\) -可积集。证明 \(x \mapsto \rho(xA, A)\) 是连续函数。
\item
  设 U 是 e 的一个邻域。证明存在一个紧集 A 和一个数 \(\varepsilon > 0\) ，使得 \(\rho(x\mathbf{A}, \mathbf{A}) < \varepsilon\) 蕴含 \(x \in U\) 。（取 A 为 e 的一个使 \(A \cdot A^{-1} \subset U\) 的邻域。）
\item
  为使 G 的子集 C 相对紧，必须且只需存在一个 \(\beta\) -可积集 A 和一个数 \(a (0 < a < 2\beta(A))\) ，使得 \(x \in C\) 蕴含 \(\rho(xA, A) \leqslant a\) 。
\end{enumerate}

\begin{enumerate}
\def\labelenumi{\arabic{enumi})}
\setcounter{enumi}{10}
\tightlist
\item
  \begin{enumerate}
  \def\labelenumii{\alph{enumii})}
  \tightlist
  \item
    设 G 为由 e 的一个紧邻域生成的局部紧群。设 \(\varphi\) 是 G 的一个非满射连续自同态，它属于 G 的（双连续）自同构所成的群在紧收敛拓扑下的闭包。那么 \(\lim_{\psi\in\mathcal{G},\psi\to\varphi}\mod\psi=0\) 。（设 K 是生成 G 的 e 的一个紧邻域。设 \(\mu\) 是 G 上的一个左 Haar 测度。若 \(\mu(\varphi(K)) > 0\) ，则 \(\varphi(K) \cdot \varphi(K)^{-1}\) 是 \(e\) 在 G 中的一个邻域，因此 \(\varphi(G)\) 是 G 的开子群；对 \(\psi \in \mathcal{G}\) 充分接近 \(\varphi\) ，有 \(\psi(K) \subset \varphi(G)\) ，从而 \(\psi(G) \neq G\) ，这是荒谬的。于是 \(\mu(\varphi(K)) = 0\) 。当 \(\psi \in \mathcal{G}\) 趋于 \(\varphi\) 时，\(\mu(\psi(K))\) 趋于 0。）
  \end{enumerate}
\end{enumerate}

\begin{enumerate}
\def\labelenumi{\alph{enumi})}
\setcounter{enumi}{1}
\tightlist
\item
  设 G 是一个自由交换群，是直和 \(G_{1} \oplus G_{2} \oplus \cdots\) ，其中每个 \(G_{i}\) 都同构于 Z。把 G 看作离散的。设 \(\varphi\) 是 G 的非满射自同态 \((x_{1}, x_{2}, \ldots) \mapsto (0, x_{1}, x_{2}, \ldots)\) 。那么 \(\varphi\) 是 G 的自同构关于紧收敛拓扑的极限，且 G 的每个自同构的模都等于 1。
\end{enumerate}

\begin{enumerate}
\def\labelenumi{\arabic{enumi})}
\setcounter{enumi}{11}
\tightlist
\item
  对 \(t > 0\) 与 \(x \in \mathbf{R}\) ，令 \(F_t(x) = te^{-\pi t^2 x^2}\) 。设 \(f \in \overline{\mathcal{H}(\mathbf{R})}\) 。证明 \(f * F_t\) 一致地趋于 \(f\) （当 \(t\) 趋于 \(+\infty\) 时）。（证明测度 \(F_t(x) dx\) 满足 §2, No.~7 的引理 4 的条件，其中 \(a = 0\) 。）
\end{enumerate}

¶ 13) 设 G 为局部紧群，连续地从左边作用在一个波兰局部紧空间 T 中。设 β 是 G 上的一个左 Haar 测度，ν 是 T 上一个拟不变正测度，χ(s,x) 是 G × T 上一个 \textgreater0 的函数，满足 Ch. VII, §1 的 Exer. 13 的条件。

对 \(f \in \mathrm{L}^p(\mathrm{T}, \nu)\) （ \(1 \leqslant p < +\infty\) ），令

\[
\left(\gamma_ {\chi , p} (s) f\right) (x) = \chi (s ^ {- 1}, x) ^ {1 / p} f (s ^ {- 1} x).
\]

\begin{enumerate}
\def\labelenumi{\alph{enumi})}
\item
  证明对每个 \(s \in G\) ，\(\pmb{\gamma}_{\chi,p}(s)\) 是 \(L^p(T, \nu)\) 的一个等距自同态。证明映射 \(s \mapsto \pmb{\gamma}_{\chi,p}(s)\) 是 \(G\) 在 \(L^p(T, \nu)\) 中的一个线性表示（按 §2, No.~5 的方式论证）。
\item
  设 \(f \in \mathcal{L}^p(\mathbf{T}, \nu)\) ，并设 \(h \in \mathcal{K}(\mathbf{G})\) 。证明函数
\end{enumerate}

\[
(x, s) \mapsto f (s ^ {- 1} x) h (s) \chi (s ^ {- 1}, x) ^ {1 / p}
\]

是 \(p\) 次方可积的（对 \(\beta \otimes \nu\) ）（先从 \(p = 1\) 的情形开始，并利用 §4, No.~1 的引理 1）。证明：若 \(q\) 是与 \(p\) 共轭的指数，且 \(g \in \mathcal{L}^q(\mathrm{T}, \nu)\) ，则 \(f(s^{-1}x)h(s)g(x)\chi(s^{-1},x)^{1/p}\) 对 \(\beta\otimes\nu\) 可积（把 \(h\) 写成 \(h_1h_2\) 的形式，其中 \(h_1,h_2\) 属于 \(\mathcal{K}(\mathrm{G})\) ）。然后由 Lebesgue-Fubini 定理推出：对 \(f\in\mathrm{L}^{p}(\mathrm{T},\nu)\) 与 \(g\in\mathrm{L}^{q}(\mathrm{T},\nu)\) ，函数

\[
s \mapsto \int g (x) (\gamma_ {\chi , p} (s) f) (x) d \nu (x)
\]

是 \(\beta\) -可测的。

\begin{enumerate}
\def\labelenumi{\alph{enumi})}
\setcounter{enumi}{2}
\item
  试利用 §4, No.~6 的 Prop. 18 的推论 2 以及 Ch. VI, §3, No.~1 的引理 1 证明：表示 \(s \mapsto \gamma_{\chi,p}(s)\) （G 在 \(L^p(T, \nu)\) 中的）对 \(1 \leqslant p < +\infty\) 是连续的。
\item
  此外假设对每个 \(s \in G\) ，函数 \(x \mapsto \chi(s^{-1}, x)\) 有界。对 \(f \in L^{p}(T, \nu)\) ，令
\end{enumerate}

\[
\left(\gamma_ {\chi} (s) f\right) (x) = \chi (s ^ {- 1}, x) f (s ^ {- 1} x).
\]

证明 \(s \mapsto \gamma_{\chi}(s)\) 是 \(G\) 在 \(L^p(T, \nu)\) 中的一个连续表示（\(1 \leqslant p < +\infty\) ）。（如同 §2, No.~5 的命题 9 的证明中那样，先证明 \(s \mapsto \gamma_{\chi}(s)\) 是 \(G\) 的一个表示，以 \(L^p(T, \nu)\) 的自同态为值。然后注意：若 \(f \in L^p(T, \nu)\) 且 \(h \in \mathcal{K}(G)\) ，则 §4, No.~1 的引理 1 表明 \(h(s)f(s^{-1}x)\chi(s^{-1}, x)\) 是 \(p\) 次方可积的（对 \(\beta \otimes \nu\) ）。证明的结尾与 c) 相同。）

\begin{enumerate}
\def\labelenumi{\arabic{enumi})}
\setcounter{enumi}{13}
\tightlist
\item
  \begin{enumerate}
  \def\labelenumii{\alph{enumii})}
  \tightlist
  \item
    设 \(G\) 为局部紧群，\(f\) 是 \(G\) 上的一个下半连续正函数，\(\mu\) 是 \(G\) 上的一个正测度。证明函数
  \end{enumerate}
\end{enumerate}

\[
x \mapsto \int^ {*} f (s ^ {- 1} x) d \mu (s) = g (x)
\]

在 G 上是下半连续的（cf.~Ch. IV, §1, No.~1, Th. 1）；为使 \(\mu\) 与 \(f\) 可卷积，必须且只需 \(g\) 对 G 上的某个左 Haar 测度可积。

\begin{enumerate}
\def\labelenumi{\alph{enumi})}
\setcounter{enumi}{1}
\tightlist
\item
  在群 \(\mathbf{G} = \mathbf{R} \times \mathbf{R}\) 上，设 \(\mu\) 是测度 \(\varepsilon_0 \otimes \lambda\) ，其中 \(\lambda\) 是 \(\mathbf{R}\) 上的 Lebesgue 测度；设 \(f(x, y) = (1 - |xy - 2|)^+\) ；证明函数
\end{enumerate}

\[
g (x, y) = \int f (x - s, y - t) d \mu (s, t)
\]

在 G 上处处有限，但既不连续，也不对 G 上的 Lebesgue 测度可积。

¶ 15) 设 G 为局部紧群，β 为 G 上的一个左 Haar 测度；以下按语言上的滥用，把 β-可积的数值函数与它们在 L¹(G,β) 中的类等同起来；对 Lᵖ(G,β) 同样滥用。

\begin{enumerate}
\def\labelenumi{\alph{enumi})}
\item
  设 \(A\) 是 Banach 空间 \(\mathrm{L}^1(\mathrm{G},\beta)\) 的一个连续自同态，使得对每个 \(s \in \mathrm{G}\) 都有 \(A(f*\varepsilon_s) = A(f)*\varepsilon_s\) （对所有 \(f \in \mathrm{L}^1(\mathrm{G},\beta)\) ）。证明：于是对每个函数 \(g \in \mathcal{K}(\mathrm{G})\) 都有 \(A(f*g) = A(f)*g\) （注意 \(s \mapsto g(s)A(f*\varepsilon_s)\) 是 \(\beta\) -可积的映射（\(\mathrm{G}\) 到 \(\mathrm{L}^1(\mathrm{G},\beta)\) 中的））；反之亦然。由此推出：存在唯一的有界测度 \(\mu\) （\(\mathrm{G}\) 上的），使得 \(A(f) = \mu *f\) 对所有 \(f \in \mathrm{L}^1(\mathrm{G},\beta)\) 成立，且 \(\| A\| = \| \mu \|\) 。（用 No.~7 的命题 19 的记号，取 \(A(f_{\mathrm{V}}*g)\) 的极限（关于比 \(\mathfrak{B}\) 的截口滤子更细的一个超滤子），并利用 \(\mathcal{M}^1(\mathrm{G})\) 的单位球对弱拓扑 \(\sigma (\mathcal{M}^1(\mathrm{G}),\overline{\mathcal{K}(\mathrm{G})})\) 的紧性。）
\item
  设 \(\mu\) 是 G 上的一个有界测度；直接证明连续自同态 \(\pmb{\gamma}(\mu): f \mapsto \mu * f\) （\(\mathrm{L}^{\infty}(\mathrm{G},\beta)\) 的）的范数等于 \(\| \mu \|\) 。（化归到 \(\mu\) 具有紧支集且关于 \(|\mu|\) 具有连续密度的情形。）由此重新推出连续自同态 \(\pmb{\gamma}(\mu): f \mapsto \mu * f\) （\(\mathrm{L}^{1}(\mathrm{G},\beta)\) 的）的范数等于 \(\| \mu \|\) 。
\item
  假设 G 是紧的且 \(\mu\) 是正的。证明对 \(1 < p < +\infty\) ，连续自同态 \(\gamma(\mu): f \mapsto \mu * f\) （\(\mathrm{L}^{p}(\mathrm{G}, \beta)\) 的）的范数等于 \(\|\mu\|\) 。
\item
  取 G 为 3 阶循环群。给出 G 上一个测度 \(\mu\) 的例子，使得自同态 \(\pmb{\gamma}(\mu)\) （\(\mathrm{L}^p (\mathrm{G},\beta)\) 的）的范数严格小于 \(\| \mu \|\) （对 \(1 < p < +\infty\) ）。
\end{enumerate}

¶ 16) 记号与约定均是 Exer. 15 的。

\begin{enumerate}
\def\labelenumi{\alph{enumi})}
\tightlist
\item
  证明：为使有界测度 \(\mu\) （\(\mathbf{G}\) 上的）满足 \(\| \mu * f \|_1 = \| f \|_1\) 对所有 \(f \in \mathrm{L}^1(\mathbf{G}, \beta)\) 都成立，必须且只需 \(\mu\) 是范数为 1 的点测度。（利用自同态 \(\gamma(|\mu|)\) （\(\mathrm{L}^1(\mathbf{G}, \beta)\) 的）的范数为 \(\| \mu \|\) 这一事实（Exer. 15），证明必然有：对每个函数 \(f \in \mathcal{K}(\mathbf{G})\) ，
\end{enumerate}

\[
\left| \int f d \mu \right| = \int | f | d | \mu |;
\]

由此先推出 \(\mu = c|\mu|\) ，其中 \(c\) 是模为 1 的常数，再推出 \(\mu\) 是一个点测度。）

\begin{enumerate}
\def\labelenumi{\alph{enumi})}
\setcounter{enumi}{1}
\tightlist
\item
  取 G 为 3 阶循环群。给出 G 上一个测度 \(\mu\) 的例子，它不是点测度，使得 \(\| \mu * f \|_2 = \| f \|_2\) 对 G 上定义的每个数值函数 \(f\) 都成立。
\end{enumerate}

\begin{enumerate}
\def\labelenumi{\arabic{enumi})}
\setcounter{enumi}{16}
\tightlist
\item
  记号与约定均是 Exer. 15 的。
\end{enumerate}

\begin{enumerate}
\def\labelenumi{\alph{enumi})}
\item
  设 \(\mu\) 是 G 上的一个有界测度；为使自同态 \(\pmb{\gamma}(\mu): f \mapsto \mu * f\) （\(\mathrm{L}^p(\mathrm{G}, \beta)\) 的）（ \(1 \leqslant p < +\infty\) ）是满射的，必须且只需存在一个 \(c_p > 0\) ，使得 \(\| \pmb{\mu}^{\vee} * g \|_q \geqslant c_p \| g \|_q\) 对所有 \(g \in \mathrm{L}^q(\mathrm{G}, \beta)\) 都成立（其中 \(\frac{1}{p} + \frac{1}{q} = 1\) ）（cf.~TVS, IV, §4, No.~2, Th. 1 的 Cor. 3 之后的注）。
\item
  若 \(\mu\) 是一个以 \(\beta\) 为基的测度，且 G 不是离散的，证明 \(\gamma(\mu)\) 对 \(1 \leqslant p \leqslant +\infty\) 永远不是满射的（对 \(p < +\infty\) ，利用 \(a\) ，化归到 \(\mu\) 关于 \(\beta\) 的密度属于 \(\mathcal{K}(\mathrm{G})\) 的情形；对 \(p = +\infty\) ，注意对 \(f \in \mathrm{L}^1(\mathrm{G}, \beta)\) 与 \(g \in \mathrm{L}^\infty(\mathrm{G}, \beta)\) ，\(f * g\) 对右一致结构是一致连续的，由此直接论证）。
\item
  若 G 是交换的，\(1 \leqslant p \leqslant 2\) ，且自同态 \(\pmb{\gamma}(\mu)\) （\(\mathrm{L}^p(\mathrm{G},\beta)\) 的）是满射的，证明它是双射的（用正则化证明：若 \(\pmb{\gamma}(\mu)\) 不是单射的，则自同态 \(\pmb{\gamma}(\check{\mu})\) （\(\mathrm{L}^q(\mathrm{G},\beta)\) 的）不是单射的：利用 Ch. IV, §6, No.~5, Prop. 4 的推论）。
\item
  取 \(\mathbf{G} = \mathbf{Z}\) 且 \(\mu = \varepsilon_1 - \varepsilon_0\) ；证明 \(\pmb{\gamma}(\mu)\) 在 \(\mathrm{L}^p (\mathbf{Z},\beta)\) （凡 \(p\neq +\infty\) 者）中都是单射的，但在 \(\mathrm{L}^{\infty}(\mathbf{Z},\beta)\) 中不是，而且对 \(p\) 的任何值都不是满射的。
\end{enumerate}

¶ 18) 记号与约定均是 Exer. 15 的。

\begin{enumerate}
\def\labelenumi{\alph{enumi})}
\item
  设 \(\mu\) 是 \(G\) 上的一个有界测度，其支集至少含两个不同的点。证明存在一个紧集 \(K\) 和一个数 \(k\)（\(0 < k < 1\)），使得 \(\| \varphi_{sK} \cdot \mu \| \leqslant k \| \mu \|\) 对所有 \(s \in G\) 都成立。（若 \(t, t'\) 是 \(\mu\) 的支集里两个不同的点，取 \(K\) 为 \(e\) 的一个充分小的邻域，使 \(sK\) 不能同时与 \(tK\) 和 \(t'K\) 相交。）
\item
  设 \(\mu\) 是一个有界测度 \(\geqslant 0\) （\(G\) 上的），其支集至少含两个不同的点。证明存在一个函数 \(f \in L^{\infty}(G, \beta)\) ，它不等价于任何 \(\geqslant 0\) 的函数，使得 \(\mu * f\) 局部几乎处处 \(\geqslant 0\) （对 \(\beta\) 而言）。（利用 \(a\) ，取 \(f\) 等于 \(-1\) （在 \(K\) 上），等于一个适当的正常数（在 \(\mathbb{C}K\) 上）。）此外，若 \(\mu\) 的支集是紧的，则存在一个具有前面诸性质且具有紧支集的函数 \(f\)。
\item
  设 \(\mu_{1},\mu_{2}\) 是 G 上两个非零的有界正测度，对卷积可交换，使得 \(\mu_{1}\) 的支集 K 是紧的且含有 \(e\) 。令 \(\mu = \mu_{1} + \mu_{2}\) 。设 V 是含有 K 的 \(e\) 的一个对称紧邻域；设 \(g\) 为这样一个函数：它等于 \(\mu * \varphi_{\mathrm{V}}\) （在 \(\mathbb{C}\mathrm{V}\) 上），等于 0（在 V 上）；证明函数 \(\mu * (g - (\mu_1 * \varphi_{\mathrm{V}}))\) 局部几乎处处 \(\geqslant 0\) （在 \(\mathbb{C}(\mathrm{V}^2)\) 中）。另一方面，证明存在一个紧集 H，使函数 \(\mu_{2} * \varphi_{\mathrm{H}}\) 几乎处处 \(\geqslant 0\) （在 \(\mathrm{V}^2\) 中）。
\item
  由 c) 推出：若 \(\mu\) 是 G 上的一个有界正测度，其支集至少含两个不同的点，则在再附加下列假设之一时，存在一个属于每个 \(L^{p}(G,\beta)\) （ \(1 \leqslant p \leqslant +\infty\) ）的函数 f，它不等价于任何 \(\geqslant 0\) 的函数，且使 \(\mu * f\) 局部几乎处处 \(\geqslant 0\) ：\(\alpha\) ) G 是交换的；\(\beta\) ) 存在一点 \(a \in G\) 使 \(\mu(\{a\}) > 0\) 。（把 \(\mu\) 适当地分解为两个可交换的测度 \(\geqslant 0\) 之和。）
\end{enumerate}

¶ 19) 记号与约定均是 Exer. 15 的。

\begin{enumerate}
\def\labelenumi{\alph{enumi})}
\item
  证明：若一个正有界测度 \(\mu\) （\(G\) 上的）使得对每个函数 \(f \in L^{p}(G, \beta)\) （ \(p\) 给定，\(1 \leqslant p < \infty\) ），关系 \(\mu * f \geqslant 0\) 局部几乎处处成立就蕴含 \(f \geqslant 0\) 局部几乎处处成立，则在下列每种情形都可断定 \(\mu\) 是一个点测度：\(\alpha)\) \(G\) 是紧的；\(\beta)\) \(G\) 是离散的；\(\gamma)\) \(G\) 是交换的（利用 Exer. 18) (*)。对 \(p = +\infty\) ，无须对 \(G\) 作补充假设，同样的结论成立。
\item
  设 \(\mu\) 是一个正有界测度（\(G\) 上的），使得：\(1^{\circ} \pmb{\gamma}(\mu)\) 是 \(L^{1}(G, \beta)\) 的一个满射自同态；\(2^{\circ}\) 对每个函数 \(f \in L^{1}(G, \beta)\) ，关系 \(\mu * f \geqslant 0\) 局部几乎处处成立蕴含 \(f \geqslant 0\) 局部几乎处处成立。证明 \(\mu\) 这时是一个点测度。（注意对每个函数 \(g \in L^{\infty}(G, \beta)\) ，关系 \(\stackrel{\vee}{\mu} * g \geqslant 0\) 局部几乎处处成立蕴含 \(g \geqslant 0\) 局部几乎处处成立。）
\end{enumerate}

\begin{enumerate}
\def\labelenumi{\arabic{enumi})}
\setcounter{enumi}{19}
\tightlist
\item
  设 \(G\) 为局部紧群，\(\beta\) 为 \(G\) 上的一个左 Haar 测度。
\end{enumerate}

\begin{enumerate}
\def\labelenumi{\alph{enumi})}
\item
  设 \(f\) 是 \(G\) 上的一个有界数值函数，对左一致结构一致连续。证明对每个有界测度 \(\mu\) （\(G\) 上的），\(\mu * f\) 都对左一致结构一致连续；此外，用 No.~7 的命题 19 的记号，\(f\) （对 \(G\) 中的一致收敛拓扑而言）是函数 \(f_V * f\) 关于 \(\mathfrak{B}\) 的截口滤子的极限。
\item
  取 \(\mathbf{G} = \mathbf{T}\) ；给出一个连续函数 \(h\) （\(\mathbf{T}\) 上的）的例子，它不是 \(f*g\) 的形式，其中 \(f\) 与 \(g\) 属于 \(\mathrm{L}^2 (\mathbf{T},\beta)\) 。（利用 Ch. IV, §6 的 Exers. 15 b) 与 16。）
\end{enumerate}

¶ 21) 设 G 为局部紧群，β 为 G 上的一个左 Haar 测度；把 L¹(G,β) 典范地等同于 M¹(G) 的一个子空间。

\begin{enumerate}
\def\labelenumi{\alph{enumi})}
\item
  用 §3 的 Exer. 11 的记号，设 A 是 \(\mathcal{M}_{\mathrm{III}}\) 的一个相对紧子集，B 是 \(\mathrm{L}^1 (\mathrm{G},\beta)\) 的一个在 \(\mathcal{M}_{\mathrm{IV}}\) 中相对紧的子集。证明 \(\mathrm{A}\times \mathrm{B}\) 上映射 \((\mu ,\nu)\mapsto \mu *\nu\) （\(\mathcal{M}_{\mathrm{III}}\times \mathcal{M}_{\mathrm{IV}}\) 到 \(\mathcal{M}_{\mathrm{IV}}\) 中的）的限制是连续的。（先证明若 A 与 B 是紧的，则 \(\mathrm{A}\times \mathrm{B}\) 在此映射下的像是紧的；利用 §3 的 Exer. 10 以及 Ch. V, §5, Exer. 15 的判别准则 \(\alpha\) )。为再证明 \((\mu ,\nu)\mapsto \mu *\nu\) 连续，利用 §3 的 Exer. 11 c)。）
\item
  设 \(\mathcal{U}_s^\infty (\mathrm{G})\) 是 \(\mathbf{G}\) 上对左一致结构一致连续的有界数值函数的集合。用 \(\mathcal{T}_{\mathrm{II}}\) 记拓扑 \(\sigma (\mathcal{M}^1 (\mathrm{G}),\mathcal{U}_s^\infty (\mathrm{G}))\) ，并用 \(\mathcal{M}_{\mathrm{II}}\) 记空间 \(\mathcal{M}^1 (\mathrm{G})\) 赋以 \(\mathcal{T}_{\mathrm{II}}\) 后所得的空间。取 \(\mathbf{G} = \mathbf{R}\) ，给出测度序列 \((\mu_n)\) （\(\mathbf{G}\) 上的）对 \(\mathcal{T}_{\mathrm{II}}\) 趋于 0 的一个例子，以及函数序列 \((f_n)\) （\(\mathrm{L}^1 (\mathrm{G},\beta)\) 中的）对 \(\mathcal{T}_{\mathrm{IV}}\) （或者等价地，对拓扑 \(\sigma (\mathrm{L}^1 (\mathrm{G},\beta),\mathrm{L}^\infty (\mathrm{G},\beta))\) ）趋于 0 的一个例子，使得序列 \((\mu_n*f_n)\) 对 \(\mathcal{T}_{\mathrm{IV}}\) 不趋于 0（取 \(f_{n}(t) = \sin nt\) 于区间 \([0,\pi ]\) 中，取 \(f_{n}(t) = 0\) 于别处）。
\item
  设 A 是 \(\mathcal{M}_{\mathrm{II}}\) 的一个相对紧子集，B 是 \(\mathrm{L}^1 (\mathrm{G},\beta)\) 的一个对范数 \(\| f\| _1\) 的拓扑相对紧的子集。证明 \(\mathbf{A}\times \mathbf{B}\) 上映射 \((\mu ,f)\mapsto \mu *\textit{f}\) （\(\mathcal{M}_{\mathrm{III}}\times \mathrm{L}^1 (\mathrm{G},\beta)\) 到 \(\mathrm{L}^1 (\mathrm{G},\beta)\) 中的）的限制是连续的（L \(^1 (\mathrm{G},\beta)\) 赋以其赋范空间拓扑）。（化归为证明：对 \(f\in \mathcal{K}(\mathrm{G})\) ，映射 \(\mu \mapsto \mu *\textit{f}\) （\(\mathcal{M}_{\mathrm{II}}\) 到 \(\mathrm{L}^1 (\mathrm{G},\beta)\) 中的）在 A 上连续。利用对 \(g\in \mathrm{L}^{\infty}(\mathrm{G},\beta)\) 与 \(f\in \mathcal{K}(\mathrm{G})\) ，\(g*\textit{f}\) 对左一致结构一致连续这一事实，先证明映射 \(\mu \mapsto \mu *\textit{f}\) （\(\mathcal{M}_{\mathrm{II}}\) 到 \(\mathcal{M}_{\mathrm{IV}}\) 中的）是连续的。然后，利用 Exer. 20 a) ，化归为证明：若 C 是 \(\mathrm{L}^1 (\mathrm{G},\beta)\) 的一个对拓扑 \(\sigma (\mathrm{L}^{1}(\mathrm{G},\beta),\mathrm{L}^{\infty}(\mathrm{G},\beta))\) 紧的子集，且 \(h\in \mathcal{K}(\mathbf{G})\) ，则映射 \(\mu \mapsto \mu *h\) （C 到 \(\mathrm{L}^1 (\mathrm{G},\beta)\) 中的）当 \(\mathrm{L}^1 (\mathrm{G},\beta)\) 赋以其赋范空间拓扑时是连续的。用与 \(a\) ) 中同样的推理证明：为此只需证明 C 在此映射下的像对赋范空间拓扑是紧的。为此，利用 Šmulian 定理、C 是挤压的（§3, Exer. 10）这一事实以及 Lebesgue 定理。）
\item
  设 A 是 \(\mathcal{M}_{\mathrm{II}}\) 的一个含有 0 的相对紧子集，B 是 \(\mathcal{M}_{\mathrm{I}}\) 的一个相对紧子集；证明对每个 \(\nu_0 \in \mathrm{B}\) ，\(\mathrm{A} \times \mathrm{B}\) 上映射 \((\mu, \nu) \mapsto \mu * \nu\) （\(\mathcal{M}_{\mathrm{II}} \times \mathcal{M}_{\mathrm{I}}\) 到 \(\mathcal{M}_{\mathrm{II}}\) 中的）的限制在点 \((0, \nu_0)\) 处连续。（利用 Exer. 20 a) 与 Exer. 21 c)。）
\item
  设 A 是 \(\mathcal{M}^1 (\mathrm{G})\) 的一个有界子集，\(f\) 是 \(\mathrm{L}^1 (\mathrm{G},\beta)\) 中的一个函数；证明映射 \(\mu \mapsto \mu *f\) （\(\mathcal{M}_{\mathrm{III}}\) 到 \(\mathcal{M}_{\mathrm{IV}}\) 中的）在 A 上的限制是连续的。
\item
  设 \(\mathcal{M}_{+}^{1}(G)\) 是 G 上有界正测度的集合。证明：若 \(\mathcal{M}_{+}^{1}(G)\) 的子集 A 对 \(T_{II}\) 是紧的，则它对 \(T_{III}\) 也是紧的。（化归为证明 A 是一个挤压集。用反证法，利用对 \(f \in \mathcal{K}(G)\) ，A 在映射 \(\mu \mapsto \mu * f\) 下的像依 c) 是一个挤压集这一事实。）
\end{enumerate}

\(g)\) 证明对 \(\mathbf{G} = \mathbf{R}\) ，\(\mathcal{M}_{+}^{1}(\mathbf{G})\) 上由 \(\mathcal{T}_{\mathrm{II}}\) 与 \(\mathcal{T}_{\mathrm{III}}\) 诱导的拓扑是不同的。

¶ 22) 记号是 §4 的 Exer. 21 与 §3 的 Exer. 11 的。

\begin{enumerate}
\def\labelenumi{\alph{enumi})}
\item
  设 \((\mu_n)\) 是 \(G\) 上的有界测度序列。假设对每个函数 \(f \in L^1(G, \beta)\) ，序列 \((\mu_n * f)\) 在 \(\mathcal{M}_I\) 中收敛于 0。证明范数序列 \((\|\mu_n\|)\) 是有界的（对映射族 \(f \mapsto \mu_n * f\) （\(L^1(G, \beta)\) 到自身的）利用 Banach-Steinhaus 定理，并利用 Exer. 15）。由此推出序列 \((\mu_n)\) 在 \(\mathcal{M}_I\) 中趋于 0（利用 Exer. 20）。
\item
  设 \((\mu_n)\) 是 \(G\) 上的有界测度序列，使得序列 \((\mu_n * f)\) 对每个函数 \(f \in L^1(G, \beta)\) 淡收敛于 0；证明序列 \((\mu_n)\) 淡收敛于 0（对每个紧子集 \(K\) （\(G\) 的），按 \(a\) ) 的方式论证序列 \((|\mu_n|(K))\) 有界）。给出一个例子（取 \(G = Z\) ），其中序列 \((\|\mu_n\|)\) 不是有界的。
\item
  设 \((\mu_n)\) 是 \(G\) 上的有界测度序列，使得对每个函数 \(f \in L^1(G, \beta)\) ，序列 \((\mu_n * f)\) 在 \(\mathcal{M}_{II}\) 中收敛于 0；证明序列 \((\mu_n)\) 在 \(\mathcal{M}_{II}\) 中收敛于 0（利用 \(a\) )）。
\end{enumerate}

\begin{enumerate}
\def\labelenumi{\arabic{enumi})}
\setcounter{enumi}{22}
\item
  设 \(G\) 为局部紧群，\(\beta\) 为 \(G\) 上的一个左 Haar 测度，\(f\) 与 \(g\) 是两个局部 \(\beta\) -可积的正函数。假设 \(f\) 与 \(g\) 可卷积，且其中一个函数在可数多个紧集之并的余集上为零。证明 \(f * g\) 局部几乎处处等于一个下半连续函数（利用 No.~5 的命题 15）。
\item
  证明 No.~5 的命题 12 与 15 中的不等式不能通过在其右端插入一个常数因子 c \textless{} 1 而加以改进。
\item
  设 G 为局部紧群，\(\beta\) 为 G 上的一个左 Haar 测度，\(\mu\) 为 G 上的一个正测度。若 A 是一个 \(\mu\) -可积集而 B 是 G 中的一个 Borel 集，则函数
\end{enumerate}

\[
u: s \mapsto \mu (\mathrm{A} \cap s \mathrm{B})
\]

在 G 中是 \(\beta\) -可测的；此外，若 \(B^{-1}\) 是 \(\beta\) -可积的，则 \(u\) 也是，且

\[
\int \mu (\mathrm{A} \cap s \mathrm{B}) d \beta (s) = \mu (\mathrm{A}) \beta (\mathrm{B} ^ {- 1}).
\]

给出一个测度 \(\mu\) 的例子，使得 \(s \mapsto \mu(A \cap sB)\) 不连续；若 \(\mu\) 是一个以 \(\beta\) 为基的测度，且 A 是 \(\mu\) -可积的而 B 是一个 Borel 集，则函数 \(s \mapsto \mu(A \cap sB)\) 在 G 上连续。

¶ 26) 设 G 为局部紧群，β 为 G 上的一个左 Haar 测度。为使 L\^{}p(G,β) 的子集 H (1 ≤ p \textless{} +∞) 是相对紧的（对 p 阶平均收敛拓扑），必须且只需下列条件满足：1° H 在 L\^{}p(G,β) 中有界；2° 对每个 ε \textgreater{} 0，存在 G 的紧子集 K，使得对每个函数 f ∈ H 有 \textbar\textbar fφ\_G-K\textbar\textbar\_p ≤ ε；3° 对每个 ε \textgreater{} 0，存在 G 中 e 的邻域 V，使得对所有 f ∈ H 与 s ∈ V 有 \textbar\textbar γ(s)f - f\textbar\textbar\_p ≤ ε。（为证明这些条件是充分的，注意：若 g ∈ K(G) 且 L 是 G 的紧子集，则集合 φ\_L · H 在映射 f ↦ g * f 下的像是 K(G) 的一个等度连续子集。）

¶ 27) 设 G 为局部紧群，β 为 G 上的一个左 Haar 测度。若 G 不约化为 e，则 L¹(G,β) 是一个（对卷积的）具有异于 0 的零因子的代数。为作出 L¹(G,β) 中两个使 f * g = 0 的非零元素 f, g，可按如下方式进行：

\(1^{\circ}\) \(\mathbf{G}\) 含有一个紧子群 \(\mathbf{H}\) （不约化为 \(e\) ），且在其中 \(\Delta(x)=1\) 的情形。取 \(f\) 为一个集合 \(\mathbf{A}\) 的特征函数，取 \(g\) 为两个特征函数之差 \(\varphi_{s\mathbf{B}}-\varphi_{\mathbf{B}}\) ，其中 \(\mathbf{A}\) 与 \(\mathbf{B}\) 适当选取。

\(2^{\circ}\) \(\mathbf{G} = \mathbf{Z}\) 的情形。证明这时可取

\[
f (n) = \frac {1}{2 n - 1} - \frac {1}{2 n + 1}
\]

对所有 \(n \in \mathbf{Z}\) ，以及 \(g(n) = f(-n)\) 。

\(3^{\circ}\) 一般情形。先证明 G 中存在一个 \(a \neq e\) 使 \(\Delta(a) = 1\) 。于是 \(a\) 生成的子群在 G 中的闭包 H 或者是一个紧子群，或者是一个同构于 Z 的子群 (GT, V, §1, Exer. 2)。第一种情形用 \(1^{\circ}\) 的结果；第二种情形取

\[
f (t) = \sum_ {n = - \infty} ^ {+ \infty} \alpha_ {n} \varphi_ {\mathrm{U} a ^ {- n}} (t), \quad g (t) = \sum_ {n = - \infty} ^ {+ \infty} \beta_ {n} \varphi_ {a ^ {n} \mathrm{U}} (t),
\]

适当选取 U 与序列 \((\alpha_{n})\) 、\((\beta_{n})\) （借助 \(2^{\circ}\) ）。

\begin{enumerate}
\def\labelenumi{\arabic{enumi})}
\setcounter{enumi}{27}
\tightlist
\item
  设 \(G_1, G_2\) 为两个局部紧群，\(\beta_1, \beta_2\) 为 \(G_1, G_2\) 上的左 Haar 测度，\(A_1\) （相应地 \(A_2\) ）为（ \(R\) 上的）拓扑代数 \(L^1(G_1, \beta_1)\) （相应地 \(L^1(G_2, \beta_2)\) ）。设 \(T\) 是 \(A_1\) 到 \(A_2\) 上的一个代数同构，使得关系 \(f \geqslant 0\) 几乎处处成立等价于 \(T(f) \geqslant 0\) 几乎处处成立。
\end{enumerate}

\begin{enumerate}
\def\labelenumi{\alph{enumi})}
\item
  证明 \(\mathbf{T}\) 是拓扑代数的一个同构。（先注意：若一个递减序列 \((f_n)\) （由 \(\mathbf{A}_1\) 的元素组成）几乎处处趋于 0，则诸 \(\mathrm{T}(f_n)\) 构成的序列也几乎处处趋于 0；由此推出：对每个序列 \((f_n)\) （\(\mathbf{A}_1\) 中有上界的），有 \(\mathrm{T}(\sup (f_n)) = \sup (\mathrm{T}(f_n))\) ，并借助 Fatou 引理得出结论。）
\item
  证明 T 可以且只可以用一种方式扩张为拓扑代数 \(\mathcal{M}^{1}(G_{1})\) 到拓扑代数 \(\mathcal{M}^{1}(G_{2})\) 上的同构，且存在 \(G_{1}\) 到 \(G_{2}\) 上的一个拓扑同构 u 与一个连续同态 \(\chi\) （\(G_{1}\) 到 \(R_{+}^{*}\) 中的），使得 \(\mathrm{T}(\varepsilon_{s}) = \chi(s)\varepsilon_{u(s)}\) （利用 Exer. 15 a) 与 Exer. 19 b)；为证明 u 与 \(\chi\) 连续，注意 T 定义了代数 \(\mathcal{L}(A_{1})\) （拓扑向量空间 \(A_{1}\) 的连续自同态所成的）到类似的代数 \(\mathcal{L}(A_{2})\) 上的一个同构，且这个同构对逐点收敛拓扑是双连续的；另一方面，注意 \(s \mapsto \delta(s^{-1})\) 是 \(G_{1}\) 到 \(\mathcal{L}(A_{1})\) 的一个乘法子群上的一个同构。
\item
  由 \(b)\) 推出：对每个 \(f \in \mathbf{A}_1\) ，
\end{enumerate}

\[
(T (f)) (t) = \chi (u ^ {- 1} (t)) f (u ^ {- 1} (t))
\]

对所有 \(t \in \mathbf{G}_2\) 成立。

回忆 (A, III, §2, Exer. 9)：存在有限群 \(\mathbf{G}_1, \mathbf{G}_2\) ，使得代数 \(\mathrm{L}^1(\mathrm{G}_1, \beta_1)\) 与 \(\mathrm{L}^1(\mathrm{G}_2, \beta_2)\) 同构而 \(\mathbf{G}_1\) 与 \(\mathbf{G}_2\) 并不同构。

\exercisesection{}

\begin{enumerate}
\def\labelenumi{\arabic{enumi})}
\tightlist
\item
  设 X 为局部紧空间，\(\mathfrak{F}\) （或 \(\mathfrak{F}(X)\) ）为 X 的闭子集的集合。对 X 的每个紧子集 K、K 的每个紧邻域 L 以及 L 的唯一一致结构的每个邻偶 V，令 \(Q(K,L,V)\) 为 \(\mathfrak{F}\) 中满足下列两个条件的元素对 (A,B) 的集合
\end{enumerate}

\[
\mathrm{A} \cap \mathrm{K} \subset \mathrm{V} (\mathrm{B} \cap \mathrm{L}) \quad \text { and } \quad \mathrm{B} \cap \mathrm{K} \subset \mathrm{V} (\mathrm{A} \cap \mathrm{L}).
\]

\begin{enumerate}
\def\labelenumi{\alph{enumi})}
\item
  证明诸集合 \(\mathbf{Q}(\mathbf{K},\mathbf{L},\mathbf{V})\) 构成 \(\mathfrak{F}\) 上一个分离一致结构的邻偶基本系。
\item
  设 \(\mathcal{U}\) 是一个与 X 的拓扑相容的一致结构。对 X 的每个紧子集 K 与 \(\mathcal{U}\) 的每个邻偶 W，令 P(K, W) 为 \(\mathfrak{F}\) 中满足下列两个条件的元素对 (A, B) 的集合
\end{enumerate}

\[
\mathrm{A} \cap \mathrm{K} \subset \mathrm{W} (\mathrm{B}) \quad \text { and } \quad \mathrm{B} \cap \mathrm{K} \subset \mathrm{W} (\mathrm{A}).
\]

证明诸集合 \(\mathbf{P}(\mathbf{K},\mathbf{W})\) 构成 \(a\) ) 中定义的一致结构的邻偶基本系。

\begin{enumerate}
\def\labelenumi{\alph{enumi})}
\setcounter{enumi}{2}
\tightlist
\item
  设 \(X'\) 是 X 的 Alexandroff 紧化，\(\omega\) 是 \(X'\) 的无穷远点；对 X 的每个闭子集 A，令 \(A' = A \cup \{\omega\}\) 。证明映射 \(A \mapsto A'\) 是 a) 中定义的一致空间 \(\mathfrak{F}(X)\) 到 \(\mathfrak{F}(X')\) 的由含 \(\omega\) 的闭子集构成的子空间上的一个同构（利用 b)）。由此推出 \(\mathfrak{F}(X)\) 是紧的（cf.~GT, II, §4, Exer. 15）。
\end{enumerate}

\begin{enumerate}
\def\labelenumi{\arabic{enumi})}
\setcounter{enumi}{1}
\item
  设 \(G\) 为局部紧群，\(\mathfrak{F}(G)\) 为 \(G\) 的闭子集的一致空间（在 No.~6 或 Exer. 1 中定义）。直接证明闭子群的集合 \(\Sigma\) （\(G\) 的）在 \(\mathfrak{F}(G)\) 中是闭的（不利用 No.~1 的命题 3）。
\item
  设 \(G\) 为局部紧群，\(\chi\) 是 \(G\) 在 \(\mathbf{R}_+^*\) 中的一个连续表示。把命题 4、5 与 6 推广到集合 \(\Gamma_{\chi}\) （诸测度 \(\alpha \in \Gamma\) 所成的），其中 \(\chi\) 在 \(H_{\alpha}\) 上的限制是 \(\alpha\) 的模。（把测度 \(\mu\) 换成 \(\chi \cdot \mu\) ，并考虑商测度 \((\chi \cdot \mu) / \overset{\vee}{\alpha}\) 。）
\item
  设 G 为一个不能由 G 的任何紧子集生成的局部紧群，并设 H 是 G 的一个使 G/H 为紧的离散子群。证明 H 不能由有限多个元素生成，但 H 的规范化 Haar 测度 \(\alpha_{0}\) 属于 \(\Gamma^{0}\) 中诸 \(\alpha\) （满足 \(\|\mu_{\alpha}\| = +\infty\) 者）所成集合的闭包（考虑 H 的由有限多个元素生成的诸子群的规范化 Haar 测度）。
\item
  设 G 是一个作为二元素子群的无穷序列 \((\mathrm{G}_{n})\) 的直和的离散群，并设 \(p_{n}\) 是 G 到 \(G_{n}\) 上的投影。令 \(\mathrm{H}_{n} = \overline{p_{n}}(e)\) ，并设 \(\alpha_{n}\) 是 \(H_{n}\) 的规范化 Haar 测度。若 \(\alpha\) 是 G 的规范化 Haar 测度，证明在 \(\Gamma_{c}^{0}\) 中序列 \((\alpha_{n})\) 收敛于 \(\alpha\) ，但 \(\|\mu_{\alpha_{n}}\| = 2\) 且 \(\|\mu_{\alpha}\| = 1\) 。
\item
  设 G 为不满足 No.~4 的条件 (L) 的局部紧群。
\end{enumerate}

\begin{enumerate}
\def\labelenumi{\alph{enumi})}
\item
  证明（用 No.~4 的记号）：N 到 D 上的典范双射不是同胚。（对 G 中 e 的每个邻域 W，令 A(W) 为含于 W 中的有限子群 \(\neq\{e\}\) 所成的集合；证明诸 A(W) 构成一个在 D 中收敛于子群 \(\{e\}\) 的滤子的基，但若 \(\alpha_{H}\) 记 H 的规范化 Haar 测度，则诸测度 \(\alpha_{H}\) 不收敛于 \(\varepsilon_{e}\) 。）
\item
  此外若 G 是可度量的，证明 G 的离散子群空间 D 中存在一个紧子集 B，它不含于 No.~3 的 Th. 1 中定义的任何集合 \(D_{U}\) 之中。
\end{enumerate}

\specialchapter{历史注记（第七章与第八章）}
